\chapter{Il livello fisico}

\begin{center}
\begin{tikzpicture}[node distance=0pt]
  \node[lay app, minimum width=3.4cm] (a) {Applicazione};
  \node[lay tra, minimum width=3.4cm, below=of a] (t) {Trasporto};
  \node[lay net, minimum width=3.4cm, below=of t] (n) {Rete};
  \node[lay link, minimum width=3.4cm, below=of n] (l) {Collegamento};
  \node[lay phy, minimum width=3.4cm, below=of l, line width=1.5pt, draw=netred] (p) {\bfseries Fisico};
  \draw[-{Stealth}, thick, netred] ([xshift=0.3cm]p.east) -- ++(1,0) node[right, lbl, text=netred] {siamo qui};
\end{tikzpicture}
\end{center}

Tutto ciò che abbiamo visto finora finiva con un frame consegnato all'adattatore. In questo capitolo vediamo che cosa ne fa
l'adattatore. Ogni altro livello aggiunge un header; questo trasforma i bit in un \textbf{segnale} (una tensione, una luce, un'onda
radio) ed è l'unico che sposta davvero qualcosa. Le domande sono quattro:

\begin{itemize}
    \item \textbf{quale mezzo}: rame, vetro o aria, e che cosa può portare ciascuno;
    \item \textbf{quale segnale}: come un bit diventa un segnale, e come il ricevente lo riporta a bit;
    \item \textbf{quanto}: ogni canale ha un tetto alla velocità in bit, e a decidere dove sta è la natura, non l'ingegneria;
    \item \textbf{quanti insieme}: posare un cavo costa, quindi un cavo deve portare molte conversazioni.
\end{itemize}

% ══════════════════════════════════════════════════════════════
\section{I mezzi trasmissivi}
% ══════════════════════════════════════════════════════════════

A ogni hop un bit viene trasmesso di nuovo, da un trasmettitore a un ricevitore, attraverso un \textbf{mezzo fisico}, e il mezzo
cambia lungo il percorso: rame fino allo switch, fibra fino al provider, radio fino al telefono in tasca.

\dfn{Mezzi guidati e non guidati}{
    \begin{itemize}
      \item Nei \textbf{mezzi guidati} l'onda è guidata lungo un mezzo solido: doppino di rame, cavo coassiale, fibra ottica.
      \item Nei \textbf{mezzi non guidati} l'onda si propaga nell'aria o nello spazio: una LAN wireless, una rete cellulare, un canale
            satellitare.
    \end{itemize}
}

Il cavo costa poco, posarlo no: la manodopera costa ordini di grandezza più del materiale, e per questo gli edifici vengono cablati
per più di quanto serva.

\begin{esempio}{Il collegamento più veloce è un camion}
Una scatola di nastri da 800 terabyte, consegnata ovunque nel Paese in 24 ore. Qual è la sua velocità in bit? 6400 terabit in 86\,400
secondi: circa \textbf{74 Gb/s}. A un'ora di strada, 1,8 Tb/s: nessun collegamento dell'edificio ci si avvicina. E si fa davvero:
Amazon sposta interi datacenter con un camion che porta 100 PB di dischi. Il problema: il primo bit arriva un giorno dopo essere
partito. Il \textbf{rate} è enorme, il \textbf{ritardo} si misura in ore. Rate e ritardo sono indipendenti, come lo erano ritardo di
trasmissione e di propagazione, e un mezzo si giudica su entrambi. Da qui in poi parliamo di mezzi che rispondono in microsecondi.
\end{esempio}

\subsection{Il doppino intrecciato}

\begin{itemize}
    \item Due fili di rame isolati, spessi circa 1 mm, \textbf{intrecciati} tra loro: il mezzo più vecchio, e ancora il più comune.
    \item Due fili paralleli sono un'ottima antenna. Intrecciati, le onde emesse dai giri successivi si annullano, e le coppie di un
          fascio smettono di disturbarsi a vicenda.
    \item Il segnale è la \textbf{differenza} tra i due fili: il rumore colpisce entrambi allo stesso modo, e la differenza gli
          sopravvive.
    \item Categorie: \textbf{Cat 5e}, quattro coppie in una guaina; \textbf{Cat 6A}, 10 Gb/s fino a 100 m; \textbf{Cat 8}, più veloce
          ma solo per circa 30 m.
    \item Ethernet a 100 Mb/s usa due coppie, una per direzione; a 1 Gb/s le usa tutte e quattro, in entrambe le direzioni
          contemporaneamente.
    \item La coppia che arriva a casa porta il telefono, e sopra di esso l'ADSL.
\end{itemize}

\subsection{Il cavo coassiale, e perfino la rete elettrica}

\begin{itemize}
    \item Un'anima di rame, un isolante, un conduttore esterno a calza, una guaina: due conduttori \textbf{concentrici} invece che
          paralleli.
    \item La schermatura compra più banda e un'immunità al rumore molto migliore di un doppino non schermato: fino a circa 6 GHz.
    \item La televisione via cavo e Internet via cavo spostano ogni segnale nella propria banda di frequenze, e molte bande
          condividono un cavo. Può essere un \textbf{mezzo condiviso}: ogni stazione collegata sente ciò che inviano le altre.
    \item \textbf{Rete elettrica} (\emph{powerline}): i dati si sovrappongono al segnale di potenza a 50 Hz. Ma quei cavi non sono mai
          stati pensati per i megahertz: attenuano, irradiano, e cambiano comportamento ogni volta che si accende un elettrodomestico.
\end{itemize}

\subsection{La luce intrappolata nel vetro}

Una sorgente di luce, una fibra di vetro, un rivelatore: un impulso di luce è un 1, il buio uno 0. Il sistema è \textbf{simplex}
(una sola direzione). La luce che incontra il bordo del nucleo con un angolo oltre quello \textbf{critico} viene riflessa all'interno,
e viaggia per chilometri quasi senza perdite.

\begin{figure}[H]
\centering
\begin{tikzpicture}[every node/.style={font=\scriptsize\sffamily}]
  \fill[netblue!8] (0,0) rectangle (9,1.6); \fill[netblue!25] (0,0.3) rectangle (9,1.3);
  \draw[black!40] (0,0) -- (9,0) (0,1.6) -- (9,1.6);
  \draw[netred, thick] (0,0.8) -- (0.8,1.3) -- (2.4,0.3) -- (4,1.3) -- (5.6,0.3) -- (7.2,1.3) -- (8.8,0.3) -- (9,0.43);
  \draw[netorange, thick] (0,0.8) -- (1.5,1.3) -- (4.5,0.3) -- (7.5,1.3) -- (9,0.8);
  \draw[netgreen!60!black, thick, -{Stealth}] (0,0.8) -- (0.4,1.3) -- (0.75,1.95);
  \node[anchor=west, text=netgreen!50!black] at (0.9,1.95) {troppo inclinato: esce dal nucleo};
  \node[anchor=west] at (9.1,1.45) {mantello}; \node[anchor=west] at (9.1,0.8) {nucleo $\approx 50\ \mu$m};
  \node[font=\footnotesize\bfseries\sffamily] at (4.5,-0.3) {multimodale: un raggio per ogni angolo};
  \begin{scope}[yshift=-2.3cm]
    \fill[netblue!8] (0,0) rectangle (9,1.2); \fill[netblue!25] (0,0.5) rectangle (9,0.7);
    \draw[black!40] (0,0) -- (9,0) (0,1.2) -- (9,1.2);
    \draw[netred, thick, -{Stealth}] (0,0.6) -- (9,0.6);
    \node[anchor=west] at (9.1,1.0) {mantello}; \node[anchor=west] at (9.1,0.6) {nucleo 8--10 $\mu$m};
    \node[font=\footnotesize\bfseries\sffamily] at (4.5,-0.3) {monomodale: dritto lungo il nucleo};
  \end{scope}
\end{tikzpicture}
\caption{Fibra multimodale e monomodale.}
\end{figure}

\begin{itemize}
    \item \textbf{Multimodale}: un nucleo largo, molti raggi a molti angoli, che arrivano in tempi diversi.
    \item \textbf{Monomodale}: un nucleo largo poche lunghezze d'onda, dove la luce va dritta; arriva 50 volte più lontano: 100 Gb/s
          per 100 km senza amplificazione.
\end{itemize}

\subsubsection*{Fibra contro rame}

\begin{itemize}
    \item \textbf{Banda}: si usano tre bande di lunghezze d'onda, attorno a 0,85, 1,30 e 1,55 $\mu$m, ciascuna larga da 25\,000 a
          30\,000 GHz.
    \item \textbf{Attenuazione}: ripetitori ogni 50 km sulle linee lunghe, contro circa 5 km sul rame.
    \item \textbf{Immunità}: niente interferenze elettromagnetiche né sbalzi di tensione; la luce non esce, quindi una fibra è difficile
          da intercettare.
    \item \textbf{Peso}: mille doppini lunghi 1 km pesano 8000 kg; due fibre portano di più e pesano 100 kg.
    \item \textbf{Contro}: si rompe se piegata troppo, va in una sola direzione (due fibre per collegamento), e le interfacce costano di
          più.
    \item Circa 100 Gb/s per canale non è il limite del vetro: è la velocità con cui sappiamo convertire tra elettricità e luce.
\end{itemize}

\begin{table}[H]
\centering
\begin{tabular}{lllll}
\toprule
\textbf{Tipo} & \textbf{Segnale} & \textbf{Distanza} & \textbf{Banda} & \textbf{Punto debole} \\
\midrule
Senza fili & radio & 30--50 m & 1,3 Gb/s (Wi-Fi 5), 9,6 (Wi-Fi 6), condivisi & sicurezza, velocità \\
Doppino & elettrico & 100 m; 30 m & 10 Gb/s (Cat 6A); 25--40 (Cat 8) & interferenze \\
Fibra ottica & luce & 100--200 km per tratta & 1--400+ Gb/s per lunghezza d'onda & fragilità \\
\bottomrule
\end{tabular}
\caption{Tre mezzi, tre compromessi.}
\end{table}

I confini si spostano ogni pochi anni: le cifre vanno lette come ordini di grandezza, non come specifiche. Quelle radio sono velocità
di picco condivise da tutta la cella: un portatile con due antenne ne vede una frazione.

% ══════════════════════════════════════════════════════════════
\section{Lo spettro elettromagnetico}
% ══════════════════════════════════════════════════════════════

Radio, microonde, infrarosso e luce sono \textbf{un unico fenomeno}, che differisce solo per la frequenza. Lunghezza d'onda $\lambda$ e
frequenza $f$ sono legate dalla velocità della luce:
$$\lambda f = c.$$
A 2,4 GHz l'onda è lunga 12,5 cm; una fibra a 1,55 $\mu$m lavora a circa 190 THz. Nel rame e nel vetro la velocità è circa $2/3$ di
$c$: i $2 \cdot 10^8$ m/s che stanno dietro a ogni ritardo di propagazione che abbiamo calcolato.

\begin{figure}[H]
\centering
\begin{tikzpicture}[every node/.style={font=\scriptsize\sffamily}]
  \fill[netorange!25] (0,0) rectangle (4.4,0.6); \node at (2.2,0.3) {radio};
  \fill[netgreen!25] (4.4,0) rectangle (7.7,0.6); \node at (6.05,0.3) {microonde};
  \fill[netred!20] (7.7,0) rectangle (11.6,0.6); \node at (9.65,0.3) {infrarosso};
  \shade[left color=red!60, right color=violet!60] (11.6,0) rectangle (12,0.6);
  \fill[netblue!20] (12,0) rectangle (13.2,0.6); \node at (12.6,0.3) {UV};
  \draw[-{Stealth}] (0,0) -- (13.6,0) node[right] {$f$ (Hz)};
  \foreach \e [count=\i from 0] in {4,6,8,10,12,14,16} \draw (\i*2.2,0) -- (\i*2.2,-0.1) node[below] {$10^{\e}$};
  \foreach \x/\t in {2.2/{radio AM}, 4.4/{radio FM}, 5.92/{Wi-Fi 2,4 GHz}, 11.31/{fibra 1,55 $\mu$m}} {
    \draw[netblue, thick] (\x,0.6) -- (\x,1.05); \node[above, text=netblue] at (\x,1.0) {\t};
  }
  \node[text=netblue, anchor=west] at (6.4,1.5) {visibile};
  \draw[netblue, -{Stealth}] (7.5,1.5) -- (11.7,0.65);
\end{tikzpicture}
\caption{Lo spettro elettromagnetico, in scala logaritmica: mezzi diversi, un solo fenomeno.}
\end{figure}

\subsection{Che cosa decide la frequenza}

\begin{itemize}
    \item Le \textbf{basse frequenze} attraversano i muri e seguono il terreno, ma portano poco: la radio AM.
    \item Le \textbf{alte frequenze} viaggiano in linea retta, rimbalzano sugli ostacoli e, sopra circa 4 GHz, vengono assorbite dalla
          pioggia.
    \item \textbf{Attenuazione di percorso} (\emph{path loss}): un'onda radio si distribuisce su una superficie che cresce, quindi la sua
          potenza cala almeno come $1/r^2$.
    \item Le copie riflesse arrivano in ritardo e fuori fase, e possono cancellare l'onda diretta: \textbf{multipath fading}.
    \item Un cavo perde la stessa frazione a ogni metro; la radio la perde ogni volta che la distanza raddoppia. La radio arriva
          lontano, e \textbf{disturba lontano}: per questo quasi tutto lo spettro è soggetto a \textbf{licenza}. Verso l'infrarosso le onde
          si comportano come la luce, e si fermano al muro.
\end{itemize}

\subsection{Le bande libere, e perché il segnale si allarga}

Alcune bande sono aperte a chiunque, le bande \textbf{ISM} (\emph{Industrial, Scientific, Medical}): 2,4 GHz quasi ovunque, 5 e ora
6 GHz per il Wi-Fi. Il prezzo dell'assenza di licenza è un tetto alla potenza di trasmissione e l'obbligo di \textbf{allargare} il
segnale su molte frequenze (\emph{spread spectrum}):

\begin{itemize}
    \item \textbf{salto di frequenza} (\emph{frequency hopping}): il trasmettitore cambia frequenza centinaia di volte al secondo, come
          Bluetooth. È difficile da disturbare, e una frequenza cattiva non lo ferma;
    \item \textbf{sequenza diretta} (\emph{direct sequence}): un codice allarga ogni bit su una banda ampia, e codici diversi possono
          condividerla: è l'idea dietro il CDMA.
\end{itemize}

Così il Wi-Fi, il mouse e il forno a microonde condividono i 2,4 GHz, e si disturbano \textbf{per progetto}.

% ══════════════════════════════════════════════════════════════
\section{Un segnale è una somma di sinusoidi}
% ══════════════════════════════════════════════════════════════

Su un filo i bit viaggiano come una grandezza fisica che cambia nel tempo, una tensione o un'intensità luminosa: una funzione
$g(t)$.

\dfn{Serie di Fourier}{
    Ogni funzione periodica ragionevolmente regolare, di periodo $T$, è una somma di seni e coseni, con $f = 1/T$:
    $$g(t) = \frac{c}{2} + \sum_{n=1}^{\infty} a_n \sin(2\pi n f t) + \sum_{n=1}^{\infty} b_n \cos(2\pi n f t).$$
    $f$ è la \textbf{frequenza fondamentale}, il termine $n$-esimo la \textbf{$n$-esima armonica}. Un messaggio finito viene trattato
    come se si ripetesse all'infinito.
}

Ciò che conta qui: un canale tratta \textbf{ciascuna armonica separatamente}, lasciandone passare alcune quasi intatte e attenuandone
altre. Ogni mezzo lascia passare le frequenze fino a un certo taglio e indebolisce le altre, e un'attenuazione disuguale è una
\textbf{distorsione}.

\begin{figure}[H]
\centering
\begin{tikzpicture}[every node/.style={font=\scriptsize\sffamily}]
  \node[font=\scriptsize\ttfamily] at (0.375,1.65) {0};
  \node[font=\scriptsize\ttfamily] at (1.125,1.65) {1};
  \node[font=\scriptsize\ttfamily] at (1.875,1.65) {1};
  \node[font=\scriptsize\ttfamily] at (2.625,1.65) {0};
  \node[font=\scriptsize\ttfamily] at (3.375,1.65) {0};
  \node[font=\scriptsize\ttfamily] at (4.125,1.65) {0};
  \node[font=\scriptsize\ttfamily] at (4.875,1.65) {1};
  \node[font=\scriptsize\ttfamily] at (5.625,1.65) {0};
  \node[font=\scriptsize\ttfamily] at (7.975,1.65) {0};
  \node[font=\scriptsize\ttfamily] at (8.725,1.65) {1};
  \node[font=\scriptsize\ttfamily] at (9.475,1.65) {1};
  \node[font=\scriptsize\ttfamily] at (10.225,1.65) {0};
  \node[font=\scriptsize\ttfamily] at (10.975,1.65) {0};
  \node[font=\scriptsize\ttfamily] at (11.725,1.65) {0};
  \node[font=\scriptsize\ttfamily] at (12.475,1.65) {1};
  \node[font=\scriptsize\ttfamily] at (13.225,1.65) {0};
  \draw[-{Stealth}, black!60] (0.00,-0.25) -- (6.30,-0.25);
  \draw[black!35, dashed] (0.000,0.000) -- (0.750,0.000) -- (0.750,1.300) -- (1.500,1.300) -- (1.500,1.300) -- (2.250,1.300) -- (2.250,0.000) -- (3.000,0.000) -- (3.000,0.000) -- (3.750,0.000) -- (3.750,0.000) -- (4.500,0.000) -- (4.500,1.300) -- (5.250,1.300) -- (5.250,0.000) -- (6.000,0.000);
  \draw[very thick, netblue] plot[smooth] coordinates {(0.000,0.609) (0.037,0.620) (0.075,0.631) (0.113,0.642) (0.150,0.653) (0.188,0.663) (0.225,0.674) (0.263,0.684) (0.300,0.693) (0.338,0.702) (0.375,0.711) (0.412,0.720) (0.450,0.728) (0.487,0.736) (0.525,0.744) (0.562,0.751) (0.600,0.758) (0.637,0.764) (0.675,0.770) (0.713,0.775) (0.750,0.780) (0.787,0.785) (0.825,0.789) (0.863,0.792) (0.900,0.795) (0.938,0.798) (0.975,0.800) (1.012,0.802) (1.050,0.803) (1.087,0.804) (1.125,0.804) (1.163,0.804) (1.200,0.803) (1.238,0.802) (1.275,0.800) (1.312,0.798) (1.350,0.795) (1.387,0.792) (1.425,0.789) (1.462,0.785) (1.500,0.780) (1.538,0.775) (1.575,0.770) (1.613,0.764) (1.650,0.758) (1.688,0.751) (1.725,0.744) (1.762,0.736) (1.800,0.728) (1.837,0.720) (1.875,0.711) (1.913,0.702) (1.950,0.693) (1.988,0.684) (2.025,0.674) (2.062,0.663) (2.100,0.653) (2.138,0.642) (2.175,0.631) (2.212,0.620) (2.250,0.609) (2.288,0.597) (2.325,0.585) (2.362,0.573) (2.400,0.561) (2.438,0.549) (2.475,0.537) (2.513,0.525) (2.550,0.512) (2.587,0.500) (2.625,0.488) (2.663,0.475) (2.700,0.463) (2.737,0.450) (2.775,0.438) (2.812,0.426) (2.850,0.414) (2.888,0.402) (2.925,0.390) (2.962,0.378) (3.000,0.366) (3.038,0.355) (3.075,0.344) (3.112,0.333) (3.150,0.322) (3.188,0.312) (3.225,0.301) (3.263,0.291) (3.300,0.282) (3.337,0.273) (3.375,0.264) (3.413,0.255) (3.450,0.247) (3.487,0.239) (3.525,0.231) (3.562,0.224) (3.600,0.217) (3.638,0.211) (3.675,0.205) (3.712,0.200) (3.750,0.195) (3.788,0.190) (3.825,0.186) (3.862,0.183) (3.900,0.180) (3.938,0.177) (3.975,0.175) (4.013,0.173) (4.050,0.172) (4.088,0.171) (4.125,0.171) (4.162,0.171) (4.200,0.172) (4.237,0.173) (4.275,0.175) (4.312,0.177) (4.350,0.180) (4.388,0.183) (4.425,0.186) (4.463,0.190) (4.500,0.195) (4.537,0.200) (4.575,0.205) (4.612,0.211) (4.650,0.217) (4.688,0.224) (4.725,0.231) (4.763,0.239) (4.800,0.247) (4.838,0.255) (4.875,0.264) (4.912,0.273) (4.950,0.282) (4.987,0.291) (5.025,0.301) (5.062,0.312) (5.100,0.322) (5.138,0.333) (5.175,0.344) (5.213,0.355) (5.250,0.366) (5.287,0.378) (5.325,0.390) (5.362,0.402) (5.400,0.414) (5.438,0.426) (5.475,0.438) (5.513,0.450) (5.550,0.463) (5.588,0.475) (5.625,0.487) (5.662,0.500) (5.700,0.512) (5.737,0.525) (5.775,0.537) (5.812,0.549) (5.850,0.561) (5.888,0.573) (5.925,0.585) (5.963,0.597) (6.000,0.609)};
  \node[anchor=west, font=\scriptsize\bfseries\sffamily] at (6.10,1.00) {1 armonica};
  \draw[-{Stealth}, black!60] (7.60,-0.25) -- (13.90,-0.25);
  \draw[black!35, dashed] (7.600,0.000) -- (8.350,0.000) -- (8.350,1.300) -- (9.100,1.300) -- (9.100,1.300) -- (9.850,1.300) -- (9.850,0.000) -- (10.600,0.000) -- (10.600,0.000) -- (11.350,0.000) -- (11.350,0.000) -- (12.100,0.000) -- (12.100,1.300) -- (12.850,1.300) -- (12.850,0.000) -- (13.600,0.000);
  \draw[very thick, netblue] plot[smooth] coordinates {(7.600,-0.012) (7.637,-0.015) (7.675,-0.014) (7.712,-0.010) (7.750,-0.001) (7.787,0.011) (7.825,0.027) (7.862,0.046) (7.900,0.069) (7.938,0.096) (7.975,0.126) (8.012,0.160) (8.050,0.196) (8.088,0.235) (8.125,0.278) (8.162,0.322) (8.200,0.369) (8.237,0.418) (8.275,0.468) (8.312,0.520) (8.350,0.573) (8.387,0.627) (8.425,0.681) (8.463,0.736) (8.500,0.790) (8.537,0.845) (8.575,0.898) (8.612,0.950) (8.650,1.001) (8.688,1.050) (8.725,1.097) (8.762,1.142) (8.800,1.184) (8.838,1.223) (8.875,1.259) (8.912,1.292) (8.950,1.322) (8.987,1.348) (9.025,1.369) (9.062,1.387) (9.100,1.401) (9.137,1.410) (9.175,1.415) (9.213,1.416) (9.250,1.412) (9.287,1.403) (9.325,1.391) (9.362,1.374) (9.400,1.352) (9.438,1.326) (9.475,1.297) (9.512,1.263) (9.550,1.225) (9.588,1.184) (9.625,1.140) (9.662,1.092) (9.700,1.042) (9.738,0.988) (9.775,0.933) (9.812,0.875) (9.850,0.816) (9.887,0.755) (9.925,0.693) (9.962,0.630) (10.000,0.566) (10.037,0.503) (10.075,0.440) (10.113,0.377) (10.150,0.315) (10.188,0.254) (10.225,0.195) (10.262,0.137) (10.300,0.082) (10.337,0.029) (10.375,-0.021) (10.412,-0.069) (10.450,-0.113) (10.488,-0.154) (10.525,-0.191) (10.562,-0.225) (10.600,-0.254) (10.637,-0.280) (10.675,-0.302) (10.712,-0.319) (10.750,-0.332) (10.787,-0.341) (10.825,-0.346) (10.863,-0.346) (10.900,-0.342) (10.938,-0.334) (10.975,-0.322) (11.012,-0.306) (11.050,-0.286) (11.087,-0.262) (11.125,-0.235) (11.162,-0.205) (11.200,-0.171) (11.238,-0.135) (11.275,-0.096) (11.312,-0.055) (11.350,-0.012) (11.387,0.033) (11.425,0.079) (11.462,0.126) (11.500,0.175) (11.537,0.223) (11.575,0.272) (11.613,0.321) (11.650,0.369) (11.688,0.417) (11.725,0.463) (11.762,0.509) (11.800,0.552) (11.837,0.594) (11.875,0.634) (11.912,0.671) (11.950,0.706) (11.988,0.738) (12.025,0.767) (12.062,0.793) (12.100,0.816) (12.137,0.835) (12.175,0.851) (12.212,0.863) (12.250,0.872) (12.287,0.877) (12.325,0.878) (12.363,0.876) (12.400,0.870) (12.438,0.861) (12.475,0.849) (12.512,0.833) (12.550,0.814) (12.587,0.792) (12.625,0.767) (12.662,0.740) (12.700,0.711) (12.738,0.679) (12.775,0.645) (12.812,0.610) (12.850,0.573) (12.887,0.535) (12.925,0.497) (12.962,0.458) (13.000,0.419) (13.037,0.379) (13.075,0.341) (13.113,0.302) (13.150,0.265) (13.188,0.229) (13.225,0.195) (13.262,0.162) (13.300,0.132) (13.337,0.104) (13.375,0.078) (13.412,0.055) (13.450,0.035) (13.488,0.018) (13.525,0.005) (13.562,-0.005) (13.600,-0.012)};
  \node[anchor=west, font=\scriptsize\bfseries\sffamily] at (13.70,1.00) {2 armoniche};
  \draw[-{Stealth}, black!60] (0.00,-2.75) -- (6.30,-2.75);
  \draw[black!35, dashed] (0.000,-2.500) -- (0.750,-2.500) -- (0.750,-1.200) -- (1.500,-1.200) -- (1.500,-1.200) -- (2.250,-1.200) -- (2.250,-2.500) -- (3.000,-2.500) -- (3.000,-2.500) -- (3.750,-2.500) -- (3.750,-2.500) -- (4.500,-2.500) -- (4.500,-1.200) -- (5.250,-1.200) -- (5.250,-2.500) -- (6.000,-2.500);
  \draw[very thick, netblue] plot[smooth] coordinates {(0.000,-2.747) (0.037,-2.728) (0.075,-2.702) (0.113,-2.670) (0.150,-2.634) (0.188,-2.593) (0.225,-2.548) (0.263,-2.501) (0.300,-2.451) (0.338,-2.400) (0.375,-2.347) (0.412,-2.294) (0.450,-2.240) (0.487,-2.187) (0.525,-2.134) (0.562,-2.081) (0.600,-2.029) (0.637,-1.978) (0.675,-1.928) (0.713,-1.878) (0.750,-1.829) (0.787,-1.781) (0.825,-1.733) (0.863,-1.685) (0.900,-1.637) (0.938,-1.590) (0.975,-1.543) (1.012,-1.495) (1.050,-1.448) (1.087,-1.401) (1.125,-1.355) (1.163,-1.310) (1.200,-1.265) (1.238,-1.222) (1.275,-1.181) (1.312,-1.142) (1.350,-1.106) (1.387,-1.073) (1.425,-1.045) (1.462,-1.021) (1.500,-1.002) (1.538,-0.988) (1.575,-0.981) (1.613,-0.980) (1.650,-0.987) (1.688,-1.000) (1.725,-1.021) (1.762,-1.049) (1.800,-1.084) (1.837,-1.127) (1.875,-1.177) (1.913,-1.233) (1.950,-1.295) (1.988,-1.363) (2.025,-1.435) (2.062,-1.512) (2.100,-1.591) (2.138,-1.672) (2.175,-1.755) (2.212,-1.838) (2.250,-1.920) (2.288,-2.000) (2.325,-2.078) (2.362,-2.151) (2.400,-2.221) (2.438,-2.285) (2.475,-2.343) (2.513,-2.396) (2.550,-2.441) (2.587,-2.480) (2.625,-2.512) (2.663,-2.537) (2.700,-2.555) (2.737,-2.567) (2.775,-2.573) (2.812,-2.573) (2.850,-2.569) (2.888,-2.560) (2.925,-2.549) (2.962,-2.535) (3.000,-2.519) (3.038,-2.502) (3.075,-2.486) (3.112,-2.471) (3.150,-2.457) (3.188,-2.445) (3.225,-2.436) (3.263,-2.430) (3.300,-2.428) (3.337,-2.429) (3.375,-2.435) (3.413,-2.443) (3.450,-2.456) (3.487,-2.471) (3.525,-2.489) (3.562,-2.508) (3.600,-2.530) (3.638,-2.551) (3.675,-2.572) (3.712,-2.592) (3.750,-2.610) (3.788,-2.624) (3.825,-2.635) (3.862,-2.641) (3.900,-2.641) (3.938,-2.635) (3.975,-2.622) (4.013,-2.603) (4.050,-2.575) (4.088,-2.541) (4.125,-2.498) (4.162,-2.449) (4.200,-2.392) (4.237,-2.329) (4.275,-2.261) (4.312,-2.187) (4.350,-2.109) (4.388,-2.029) (4.425,-1.947) (4.463,-1.864) (4.500,-1.782) (4.537,-1.702) (4.575,-1.625) (4.612,-1.553) (4.650,-1.487) (4.688,-1.427) (4.725,-1.376) (4.763,-1.333) (4.800,-1.300) (4.838,-1.277) (4.875,-1.264) (4.912,-1.263) (4.950,-1.272) (4.987,-1.292) (5.025,-1.323) (5.062,-1.363) (5.100,-1.414) (5.138,-1.473) (5.175,-1.539) (5.213,-1.612) (5.250,-1.691) (5.287,-1.775) (5.325,-1.861) (5.362,-1.949) (5.400,-2.038) (5.438,-2.125) (5.475,-2.211) (5.513,-2.294) (5.550,-2.372) (5.588,-2.445) (5.625,-2.512) (5.662,-2.572) (5.700,-2.624) (5.737,-2.669) (5.775,-2.705) (5.812,-2.733) (5.850,-2.752) (5.888,-2.763) (5.925,-2.765) (5.963,-2.760) (6.000,-2.747)};
  \node[anchor=west, font=\scriptsize\bfseries\sffamily] at (6.10,-1.50) {4 armoniche};
  \draw[-{Stealth}, black!60] (7.60,-2.75) -- (13.90,-2.75);
  \draw[black!35, dashed] (7.600,-2.500) -- (8.350,-2.500) -- (8.350,-1.200) -- (9.100,-1.200) -- (9.100,-1.200) -- (9.850,-1.200) -- (9.850,-2.500) -- (10.600,-2.500) -- (10.600,-2.500) -- (11.350,-2.500) -- (11.350,-2.500) -- (12.100,-2.500) -- (12.100,-1.200) -- (12.850,-1.200) -- (12.850,-2.500) -- (13.600,-2.500);
  \draw[very thick, netblue] plot[smooth] coordinates {(7.600,-2.417) (7.637,-2.421) (7.675,-2.434) (7.712,-2.454) (7.750,-2.480) (7.787,-2.510) (7.825,-2.541) (7.862,-2.570) (7.900,-2.595) (7.938,-2.612) (7.975,-2.618) (8.012,-2.611) (8.050,-2.589) (8.088,-2.550) (8.125,-2.494) (8.162,-2.421) (8.200,-2.332) (8.237,-2.228) (8.275,-2.113) (8.312,-1.989) (8.350,-1.861) (8.387,-1.731) (8.425,-1.604) (8.463,-1.484) (8.500,-1.374) (8.537,-1.279) (8.575,-1.199) (8.612,-1.137) (8.650,-1.093) (8.688,-1.068) (8.725,-1.060) (8.762,-1.067) (8.800,-1.087) (8.838,-1.117) (8.875,-1.153) (8.912,-1.191) (8.950,-1.228) (8.987,-1.261) (9.025,-1.287) (9.062,-1.303) (9.100,-1.309) (9.137,-1.303) (9.175,-1.287) (9.213,-1.261) (9.250,-1.228) (9.287,-1.191) (9.325,-1.152) (9.362,-1.116) (9.400,-1.085) (9.438,-1.065) (9.475,-1.058) (9.512,-1.066) (9.550,-1.092) (9.588,-1.136) (9.625,-1.199) (9.662,-1.279) (9.700,-1.376) (9.738,-1.487) (9.775,-1.607) (9.812,-1.735) (9.850,-1.865) (9.887,-1.993) (9.925,-2.117) (9.962,-2.231) (10.000,-2.332) (10.037,-2.420) (10.075,-2.491) (10.113,-2.545) (10.150,-2.582) (10.188,-2.603) (10.225,-2.610) (10.262,-2.604) (10.300,-2.588) (10.337,-2.566) (10.375,-2.540) (10.412,-2.513) (10.450,-2.488) (10.488,-2.466) (10.525,-2.449) (10.562,-2.439) (10.600,-2.436) (10.637,-2.439) (10.675,-2.448) (10.712,-2.462) (10.750,-2.479) (10.787,-2.497) (10.825,-2.515) (10.863,-2.530) (10.900,-2.543) (10.938,-2.551) (10.975,-2.553) (11.012,-2.551) (11.050,-2.543) (11.087,-2.531) (11.125,-2.515) (11.162,-2.498) (11.200,-2.481) (11.238,-2.465) (11.275,-2.452) (11.312,-2.443) (11.350,-2.440) (11.387,-2.443) (11.425,-2.453) (11.462,-2.468) (11.500,-2.488) (11.537,-2.511) (11.575,-2.536) (11.613,-2.560) (11.650,-2.580) (11.688,-2.594) (11.725,-2.599) (11.762,-2.593) (11.800,-2.574) (11.837,-2.540) (11.875,-2.490) (11.912,-2.424) (11.950,-2.342) (11.988,-2.245) (12.025,-2.136) (12.062,-2.016) (12.100,-1.889) (12.137,-1.757) (12.175,-1.625) (12.212,-1.497) (12.250,-1.377) (12.287,-1.268) (12.325,-1.173) (12.363,-1.096) (12.400,-1.039) (12.438,-1.005) (12.475,-0.993) (12.512,-1.005) (12.550,-1.039) (12.587,-1.096) (12.625,-1.172) (12.662,-1.266) (12.700,-1.375) (12.738,-1.495) (12.775,-1.622) (12.812,-1.753) (12.850,-1.884) (12.887,-2.012) (12.925,-2.133) (12.962,-2.243) (13.000,-2.342) (13.037,-2.426) (13.075,-2.494) (13.113,-2.546) (13.150,-2.582) (13.188,-2.603) (13.225,-2.610) (13.262,-2.604) (13.300,-2.588) (13.337,-2.565) (13.375,-2.537) (13.412,-2.507) (13.450,-2.479) (13.488,-2.453) (13.525,-2.433) (13.562,-2.421) (13.600,-2.417)};
  \node[anchor=west, font=\scriptsize\bfseries\sffamily] at (13.70,-1.50) {8 armoniche};
\end{tikzpicture}
\caption{La lettera b (\texttt{01100010}) e ciò che arriva se il canale lascia passare solo le prime 1, 2, 4 o 8 armoniche
(tratteggiato il segnale originale).}
\end{figure}

Con 8 armoniche i bit si leggono; di più renderebbero solo il quadrato più bello.

\subsection{Banda più stretta, meno bit}

Una linea telefonica per la voce taglia poco sopra i 3000 Hz. Inviando 8 bit a $R$ b/s, il periodo è $8/R$ secondi e la prima
armonica sta a $R/8$ Hz: la linea lascia passare circa $3000/(R/8) = 24\,000/R$ armoniche.

\begin{table}[H]
\centering
\begin{tabular}{rrr}
\toprule
\textbf{b/s} & \textbf{prima armonica} & \textbf{armoniche passate} \\
\midrule
300 & 37,5 Hz & 80 \\
1200 & 150 Hz & 20 \\
2400 & 300 Hz & 10 \\
4800 & 600 Hz & 5 \\
9600 & 1200 Hz & 2 \\
19\,200 & 2400 Hz & 1 \\
38\,400 & 4800 Hz & 0 \\
\bottomrule
\end{tabular}
\caption{Su una linea da 3000 Hz, più veloce si va, meno armoniche passano.}
\end{table}

A 9600 b/s il segnale somiglia al pannello a due armoniche; a 38\,400 non passa nulla. Limitare la banda limita la velocità in bit,
\textbf{anche su un canale perfetto}. Più livelli aiuteranno, entro certi limiti.

\subsection{Banda in hertz, banda in bit al secondo}

All'inizio del corso la \emph{banda} era una velocità; avevamo detto che quaggiù sarebbe stata anche una \textbf{larghezza}.

\begin{itemize}
    \item \textbf{Banda analogica}: l'intervallo di frequenze che un mezzo lascia passare senza forte attenuazione, in \textbf{Hz}. È una
          proprietà del filo: materiale, spessore, lunghezza. Si indica fino al punto in cui la potenza ricevuta si è dimezzata, e
          dei filtri la restringono di proposito: un canale 802.11 è largo 20 MHz o più.
    \item \textbf{Banda digitale}: la velocità massima in bit, in \textbf{b/s}: ciò che si ottiene da quella analogica. Il tasso di cambio
          lo danno Nyquist e Shannon.
    \item \textbf{Banda base} (\emph{baseband}): da 0 fino a una certa frequenza, come sulla maggior parte dei fili. \textbf{Banda
          passante} (\emph{passband}): spostata in una banda più alta, come tutta la radio.
\end{itemize}

A portare informazione è la \textbf{larghezza} della banda, non la sua posizione.

% ══════════════════════════════════════════════════════════════
\section{I due tetti: Nyquist e Shannon}
% ══════════════════════════════════════════════════════════════

\dfn{Teorema di Nyquist (1924)}{
    Un canale \textbf{senza rumore} di banda $B$ Hz, che porta un segnale con $V$ livelli discreti, non può superare
    $$R_{\max} = 2B \log_2 V \ \text{b/s}.$$
}

\begin{itemize}
    \item Dietro c'è il campionamento: un segnale filtrato a $B$ Hz si ricostruisce completamente da $2B$ campioni al secondo, e
          campionare più in fretta non trova nulla di nuovo.
    \item Un segnale binario su un canale da 3 kHz: al massimo 6000 b/s, per quanto buona sia l'elettronica.
    \item La formula premia i livelli: 4 raddoppiano la velocità, 16 la quadruplicano, e niente nella formula dice quando fermarsi.
          Ma vale per un canale perfetto, e nessun canale è perfetto.
\end{itemize}

\subsection{Segnale, rumore e decibel}

Ogni canale porta \textbf{rumore termico}, dovuto al moto delle molecole al suo interno. Ciò che conta è il suo rapporto con il
segnale, il \textbf{rapporto segnale/rumore} $S/N$. Il rapporto copre molti ordini di grandezza, quindi si indica su scala
logaritmica, in \textbf{decibel}: $10 \log_{10}(S/N)$.

\begin{table}[H]
\centering
\begin{tabular}{lccccc}
\toprule
$S/N$ & 2 & 10 & 100 & 1000 & 10\,000 \\
dB & 3 & 10 & 20 & 30 & 40 \\
\bottomrule
\end{tabular}
\caption{Rapporti e decibel.}
\end{table}

\begin{itemize}
    \item 3 dB sono un fattore due: il punto di metà potenza a cui si indica la banda analogica.
    \item Le perdite in dB semplicemente \textbf{si sommano}: un doppino perde circa 20 dB ogni 100 m, la radio circa 6 dB ogni volta
          che la distanza raddoppia.
\end{itemize}

\dfn{Teorema di Shannon (1948)}{
    Un canale di banda $B$ Hz, con rumore termico e rapporto segnale/rumore $S/N$, non può superare
    $$C = B \log_2 (1 + S/N) \ \text{b/s},$$
    qualunque sia il numero di livelli e qualunque sia il campionamento.
}

\begin{esempio}{Il tetto dell'ADSL}
Circa 1 MHz di linea e 40 dB su un doppino corto: $10^6 \cdot \log_2 10\,001 \approx 13$ Mb/s. Lo standard originale prometteva 12.
Oltre ci sono solo due strade: \textbf{più banda} o \textbf{meno rumore} (ADSL2+ ha raddoppiato la banda). Un sistema che dichiara di
più va messo insieme alle macchine del moto perpetuo.
\end{esempio}

\subsection{Perché non usare semplicemente più livelli?}

Nyquist dice che ogni raddoppio dei livelli aggiunge un bit per simbolo. Perché non mille livelli? Per il \textbf{rumore}: livelli più
vicini del rumore non si distinguono più. Uguagliando la velocità di Nyquist alla capacità di Shannon,
$$2B \log_2 V = B \log_2 (1 + S/N) \quad\Longrightarrow\quad V = \sqrt{1 + S/N}.$$
A 40 dB sono circa 100 livelli; a 20 dB, solo 10. I due teoremi rispondono a domande diverse: Nyquist \textbf{quanto in fretta}
segnalare, Shannon \textbf{quanti bit} ogni segnale può portare. Tutto ciò che segue, codici di linea e costellazioni, è un modo di
spendere quel budget.

% ══════════════════════════════════════════════════════════════
\section{La trasmissione in banda base: i codici di linea}
% ══════════════════════════════════════════════════════════════

In banda base un bit diventa direttamente un \textbf{livello}: una tensione per l'1, un'altra per lo 0, oppure luce e buio su una fibra.
È la NRZ che abbiamo incontrato cercando i confini dei frame.

\begin{itemize}
    \item Il ricevente campiona a intervalli regolari e prende ogni campione per il livello più vicino: ciò che arriva è attenuato e
          rumoroso.
    \item La NRZ può cambiare a ogni bit, quindi a $R$ b/s le servono almeno $R/2$ Hz: Nyquist con due livelli.
    \item Quattro livelli portano 2 bit per simbolo: la velocità di simbolo si dimezza, e con essa la banda necessaria.
    \item \textbf{Baud} è il vecchio nome della velocità di simbolo. Velocità in bit $=$ velocità di simbolo $\times$ bit per simbolo, e
          le due vengono spesso confuse.
\end{itemize}

La NRZ da sola si usa raramente: i \textbf{codici di linea} la rimodellano.

\begin{figure}[H]
\centering
\begin{tikzpicture}[every node/.style={font=\scriptsize\sffamily}]
  \draw[black!20, dashed] (0.00,0.75) -- (0.00,-3.95);
  \draw[black!20, dashed] (0.80,0.75) -- (0.80,-3.95);
  \draw[black!20, dashed] (1.60,0.75) -- (1.60,-3.95);
  \draw[black!20, dashed] (2.40,0.75) -- (2.40,-3.95);
  \draw[black!20, dashed] (3.20,0.75) -- (3.20,-3.95);
  \draw[black!20, dashed] (4.00,0.75) -- (4.00,-3.95);
  \draw[black!20, dashed] (4.80,0.75) -- (4.80,-3.95);
  \draw[black!20, dashed] (5.60,0.75) -- (5.60,-3.95);
  \draw[black!20, dashed] (6.40,0.75) -- (6.40,-3.95);
  \draw[black!20, dashed] (7.20,0.75) -- (7.20,-3.95);
  \draw[black!20, dashed] (8.00,0.75) -- (8.00,-3.95);
  \draw[black!20, dashed] (8.80,0.75) -- (8.80,-3.95);
  \node[font=\scriptsize\ttfamily] at (0.40,0.6) {1};
  \node[font=\scriptsize\ttfamily] at (1.20,0.6) {0};
  \node[font=\scriptsize\ttfamily] at (2.00,0.6) {0};
  \node[font=\scriptsize\ttfamily] at (2.80,0.6) {0};
  \node[font=\scriptsize\ttfamily] at (3.60,0.6) {0};
  \node[font=\scriptsize\ttfamily] at (4.40,0.6) {1};
  \node[font=\scriptsize\ttfamily] at (5.20,0.6) {0};
  \node[font=\scriptsize\ttfamily] at (6.00,0.6) {1};
  \node[font=\scriptsize\ttfamily] at (6.80,0.6) {1};
  \node[font=\scriptsize\ttfamily] at (7.60,0.6) {1};
  \node[font=\scriptsize\ttfamily] at (8.40,0.6) {1};
  \node[anchor=east, font=\scriptsize\bfseries\sffamily] at (-0.15,0.00) {NRZ};
  \draw[black!25] (0,0.00) -- (8.80,0.00);
  \draw[thick, netblue] (0.000,0.320) -- (0.400,0.320) -- (0.400,0.320) -- (0.800,0.320) -- (0.800,-0.320) -- (1.200,-0.320) -- (1.200,-0.320) -- (1.600,-0.320) -- (1.600,-0.320) -- (2.000,-0.320) -- (2.000,-0.320) -- (2.400,-0.320) -- (2.400,-0.320) -- (2.800,-0.320) -- (2.800,-0.320) -- (3.200,-0.320) -- (3.200,-0.320) -- (3.600,-0.320) -- (3.600,-0.320) -- (4.000,-0.320) -- (4.000,0.320) -- (4.400,0.320) -- (4.400,0.320) -- (4.800,0.320) -- (4.800,-0.320) -- (5.200,-0.320) -- (5.200,-0.320) -- (5.600,-0.320) -- (5.600,0.320) -- (6.000,0.320) -- (6.000,0.320) -- (6.400,0.320) -- (6.400,0.320) -- (6.800,0.320) -- (6.800,0.320) -- (7.200,0.320) -- (7.200,0.320) -- (7.600,0.320) -- (7.600,0.320) -- (8.000,0.320) -- (8.000,0.320) -- (8.400,0.320) -- (8.400,0.320) -- (8.800,0.320);
  \node[anchor=east, font=\scriptsize\bfseries\sffamily] at (-0.15,-1.15) {NRZI};
  \draw[black!25] (0,-1.15) -- (8.80,-1.15);
  \draw[thick, netblue] (0.000,-0.830) -- (0.400,-0.830) -- (0.400,-0.830) -- (0.800,-0.830) -- (0.800,-0.830) -- (1.200,-0.830) -- (1.200,-0.830) -- (1.600,-0.830) -- (1.600,-0.830) -- (2.000,-0.830) -- (2.000,-0.830) -- (2.400,-0.830) -- (2.400,-0.830) -- (2.800,-0.830) -- (2.800,-0.830) -- (3.200,-0.830) -- (3.200,-0.830) -- (3.600,-0.830) -- (3.600,-0.830) -- (4.000,-0.830) -- (4.000,-1.470) -- (4.400,-1.470) -- (4.400,-1.470) -- (4.800,-1.470) -- (4.800,-1.470) -- (5.200,-1.470) -- (5.200,-1.470) -- (5.600,-1.470) -- (5.600,-0.830) -- (6.000,-0.830) -- (6.000,-0.830) -- (6.400,-0.830) -- (6.400,-1.470) -- (6.800,-1.470) -- (6.800,-1.470) -- (7.200,-1.470) -- (7.200,-0.830) -- (7.600,-0.830) -- (7.600,-0.830) -- (8.000,-0.830) -- (8.000,-1.470) -- (8.400,-1.470) -- (8.400,-1.470) -- (8.800,-1.470);
  \node[anchor=east, font=\scriptsize\bfseries\sffamily] at (-0.15,-2.30) {Manchester};
  \draw[black!25] (0,-2.30) -- (8.80,-2.30);
  \draw[thick, netblue] (0.000,-2.620) -- (0.400,-2.620) -- (0.400,-1.980) -- (0.800,-1.980) -- (0.800,-1.980) -- (1.200,-1.980) -- (1.200,-2.620) -- (1.600,-2.620) -- (1.600,-1.980) -- (2.000,-1.980) -- (2.000,-2.620) -- (2.400,-2.620) -- (2.400,-1.980) -- (2.800,-1.980) -- (2.800,-2.620) -- (3.200,-2.620) -- (3.200,-1.980) -- (3.600,-1.980) -- (3.600,-2.620) -- (4.000,-2.620) -- (4.000,-2.620) -- (4.400,-2.620) -- (4.400,-1.980) -- (4.800,-1.980) -- (4.800,-1.980) -- (5.200,-1.980) -- (5.200,-2.620) -- (5.600,-2.620) -- (5.600,-2.620) -- (6.000,-2.620) -- (6.000,-1.980) -- (6.400,-1.980) -- (6.400,-2.620) -- (6.800,-2.620) -- (6.800,-1.980) -- (7.200,-1.980) -- (7.200,-2.620) -- (7.600,-2.620) -- (7.600,-1.980) -- (8.000,-1.980) -- (8.000,-2.620) -- (8.400,-2.620) -- (8.400,-1.980) -- (8.800,-1.980);
  \node[anchor=east, font=\scriptsize\bfseries\sffamily] at (-0.15,-3.45) {AMI};
  \draw[black!25] (0,-3.45) -- (8.80,-3.45);
  \draw[thick, netblue] (0.000,-3.130) -- (0.400,-3.130) -- (0.400,-3.130) -- (0.800,-3.130) -- (0.800,-3.450) -- (1.200,-3.450) -- (1.200,-3.450) -- (1.600,-3.450) -- (1.600,-3.450) -- (2.000,-3.450) -- (2.000,-3.450) -- (2.400,-3.450) -- (2.400,-3.450) -- (2.800,-3.450) -- (2.800,-3.450) -- (3.200,-3.450) -- (3.200,-3.450) -- (3.600,-3.450) -- (3.600,-3.450) -- (4.000,-3.450) -- (4.000,-3.770) -- (4.400,-3.770) -- (4.400,-3.770) -- (4.800,-3.770) -- (4.800,-3.450) -- (5.200,-3.450) -- (5.200,-3.450) -- (5.600,-3.450) -- (5.600,-3.130) -- (6.000,-3.130) -- (6.000,-3.130) -- (6.400,-3.130) -- (6.400,-3.770) -- (6.800,-3.770) -- (6.800,-3.770) -- (7.200,-3.770) -- (7.200,-3.130) -- (7.600,-3.130) -- (7.600,-3.130) -- (8.000,-3.130) -- (8.000,-3.770) -- (8.400,-3.770) -- (8.400,-3.770) -- (8.800,-3.770);
\end{tikzpicture}
\caption{Gli stessi undici bit in quattro codici di linea. NRZ: il livello è il bit. NRZI: un 1 è una transizione, uno 0 nessuna.
Manchester: una transizione a metà di ogni bit, e la sua direzione è il bit (qui salita per 1). AMI: uno 0 è zero, gli 1 si alternano.}
\end{figure}

\subsection{Tenere l'orologio del ricevente}

\begin{itemize}
    \item Con la NRZ una sequenza di bit uguali è una linea piatta: quindici zeri sembrano sedici, a meno che i due orologi non
          concordino entro una frazione di bit. Un filo separato per l'orologio risolverebbe, ma un secondo filo è meglio speso per i
          dati.
    \item \textbf{Manchester}: l'orologio messo in XOR con i dati, una transizione in ogni bit su cui il ricevente si aggancia. Il prezzo
          è il \textbf{doppio della banda}. È l'Ethernet classico.
    \item \textbf{NRZI}: una transizione per un valore, nessuna per l'altro. Le sequenze del primo valore sono al sicuro, quelle del
          secondo no.
    \item \textbf{USB} usa la NRZI con lo 0 come transizione, quindi il pericolo è una lunga sequenza di 1, interrotta dal \textbf{bit
          stuffing} visto con il framing: uno 0 forzato dopo sei 1.
\end{itemize}

\subsection{4B/5B}

\begin{minipage}{0.42\linewidth}
\centering\small
\begin{tabular}{cc|cc}
\toprule
\textbf{4B} & \textbf{5B} & \textbf{4B} & \textbf{5B} \\
\midrule
0000 & 11110 & 1000 & 10010 \\
0001 & 01001 & 1001 & 10011 \\
0010 & 10100 & 1010 & 10110 \\
0011 & 10101 & 1011 & 10111 \\
0100 & 01010 & 1100 & 11010 \\
0101 & 01011 & 1101 & 11011 \\
0110 & 01110 & 1110 & 11100 \\
0111 & 01111 & 1111 & 11101 \\
\bottomrule
\end{tabular}
\end{minipage}\hfill
\begin{minipage}{0.55\linewidth}
\begin{itemize}
    \item Ogni 4 bit di dati partono come 5, da una tabella fissa scelta in modo che non ci siano mai più di tre 0 di seguito.
    \item Inviati con un codice che marca ogni 1 con una transizione (NRZI), la linea non resta mai piatta a lungo.
    \item Il costo è del \textbf{25\%}, contro il 100\% di Manchester.
    \item Dei 32 simboli ne avanzano 16: uno significa linea inattiva, un altro inizio di frame: sono le \textbf{violazioni della
          codifica} che delimitavano i frame.
\end{itemize}
\end{minipage}

\subsection{Lo scrambling, e il pacchetto killer}

Lo \textbf{scrambling} fa lo XOR dei dati con una sequenza pseudocasuale; il ricevente rifà lo XOR con la stessa sequenza e riottiene
i dati. Nessun overhead, e l'energia si distribuisce su tutte le frequenze invece di concentrarsi su poche.

\warning{Il pacchetto killer}{
    E se i dati fossero proprio uguali alla sequenza dello scrambler? Lo XOR dà tutti zeri: una linea piatta, e un orologio perso. Le
    prime realizzazioni di IP su SONET usavano una sequenza corta e prevedibile, e un utente poteva inviare \textbf{pacchetti killer} che
    facevano cadere il collegamento. Lo scrambling rende le lunghe sequenze improbabili, non impossibili: è sicuro solo con una
    sequenza lunga che nessuno possa prevedere.
}

\subsection{Segnali bilanciati}

Un segnale con tanta tensione positiva quanta negativa, anche su tratti brevi, non ha \textbf{componente continua} (DC).

\begin{itemize}
    \item Conta: il cavo coassiale e le linee accoppiate con trasformatori filtrano via la continua, e l'energia messa lì è sprecata.
    \item La media di un segnale bilanciato è una soglia di decisione pronta; uno sbilanciato deriva, e i simboli vengono letti male.
    \item \textbf{AMI}: uno 0 è zero e gli 1 successivi si alternano tra $+$ e $-$: bilanciamento comprato con un terzo livello.
    \item \textbf{8B/10B}: 8 bit come 10, scegliendo ogni simbolo o il suo complemento per annullare la disparità accumulata: mai più di
          2 di sbilanciamento, mai più di cinque bit uguali di fila.
\end{itemize}

\begin{table}[H]
\centering
\begin{tabular}{llll}
\toprule
\textbf{Ethernet} & \textbf{Mezzo} & \textbf{Codice di linea} & \textbf{Sul filo} \\
\midrule
10 Mb/s, classico & coassiale, doppino & Manchester & 20 MHz: due per bit \\
100BASE-TX & 2 coppie, Cat 5 & 4B/5B & 125 MHz \\
1000BASE-X & fibra & 8B/10B & 1,25 Gbaud \\
1000BASE-T & 4 coppie, Cat 5 & 5 livelli, 2 bit & 125 Msimboli/s per coppia \\
10 Gb/s, fibra & fibra & scrambling + 64B/66B & 3\% di overhead \\
\bottomrule
\end{tabular}
\caption{Ogni velocità di Ethernet, la sua codifica.}
\end{table}

1000BASE-T: 4 coppie $\times$ 125 Msimboli/s $\times$ 2 bit $=$ 1 Gb/s, con ogni coppia che trasmette e riceve contemporaneamente. Il
formato del frame non è mai cambiato; a ogni generazione è cambiato il modo in cui i bit diventano un segnale.

% ══════════════════════════════════════════════════════════════
\section{La trasmissione in banda passante: la modulazione}
% ══════════════════════════════════════════════════════════════

Un segnale in banda base a bassa frequenza (una voce, un'onda quadra di bit) viene trasportato attorno a una \textbf{frequenza
portante} più alta: l'informazione sta nell'\textbf{inviluppo}. Perché?

\begin{itemize}
    \item Un'antenna deve essere grande una frazione della lunghezza d'onda: a bassa frequenza servirebbero antenne enormi.
    \item Il regolatore assegna una banda, e bisogna starci dentro.
    \item Segnali diversi possono condividere un filo in bande diverse.
\end{itemize}

Una linea telefonica per la voce lascia passare solo 300--3400 Hz: qualunque cosa la attraversi, dati compresi, deve vivere in
quella finestra.

\begin{figure}[H]
\centering
\begin{tikzpicture}[every node/.style={font=\scriptsize\sffamily}]
  \draw[black!20, dashed] (0.00,0.6) -- (0.00,-3.45);
  \draw[black!20, dashed] (0.70,0.6) -- (0.70,-3.45);
  \draw[black!20, dashed] (1.40,0.6) -- (1.40,-3.45);
  \draw[black!20, dashed] (2.10,0.6) -- (2.10,-3.45);
  \draw[black!20, dashed] (2.80,0.6) -- (2.80,-3.45);
  \draw[black!20, dashed] (3.50,0.6) -- (3.50,-3.45);
  \draw[black!20, dashed] (4.20,0.6) -- (4.20,-3.45);
  \draw[black!20, dashed] (4.90,0.6) -- (4.90,-3.45);
  \draw[black!20, dashed] (5.60,0.6) -- (5.60,-3.45);
  \draw[black!20, dashed] (6.30,0.6) -- (6.30,-3.45);
  \draw[black!20, dashed] (7.00,0.6) -- (7.00,-3.45);
  \draw[black!20, dashed] (7.70,0.6) -- (7.70,-3.45);
  \draw[black!20, dashed] (8.40,0.6) -- (8.40,-3.45);
  \draw[black!20, dashed] (9.10,0.6) -- (9.10,-3.45);
  \node[font=\scriptsize\ttfamily] at (0.35,0.55) {0};
  \node[font=\scriptsize\ttfamily] at (1.05,0.55) {1};
  \node[font=\scriptsize\ttfamily] at (1.75,0.55) {0};
  \node[font=\scriptsize\ttfamily] at (2.45,0.55) {1};
  \node[font=\scriptsize\ttfamily] at (3.15,0.55) {1};
  \node[font=\scriptsize\ttfamily] at (3.85,0.55) {0};
  \node[font=\scriptsize\ttfamily] at (4.55,0.55) {0};
  \node[font=\scriptsize\ttfamily] at (5.25,0.55) {1};
  \node[font=\scriptsize\ttfamily] at (5.95,0.55) {0};
  \node[font=\scriptsize\ttfamily] at (6.65,0.55) {0};
  \node[font=\scriptsize\ttfamily] at (7.35,0.55) {1};
  \node[font=\scriptsize\ttfamily] at (8.05,0.55) {0};
  \node[font=\scriptsize\ttfamily] at (8.75,0.55) {0};
  \node[anchor=east, font=\scriptsize\bfseries\sffamily] at (-0.15,0.00) {bit};
  \draw[black!25] (0,0.00) -- (9.10,0.00);
  \draw[thick, black] (0.000,-0.300) -- (0.700,-0.300) -- (0.700,0.300) -- (1.400,0.300) -- (1.400,-0.300) -- (2.100,-0.300) -- (2.100,0.300) -- (2.800,0.300) -- (2.800,0.300) -- (3.500,0.300) -- (3.500,-0.300) -- (4.200,-0.300) -- (4.200,-0.300) -- (4.900,-0.300) -- (4.900,0.300) -- (5.600,0.300) -- (5.600,-0.300) -- (6.300,-0.300) -- (6.300,-0.300) -- (7.000,-0.300) -- (7.000,0.300) -- (7.700,0.300) -- (7.700,-0.300) -- (8.400,-0.300) -- (8.400,-0.300) -- (9.100,-0.300);
  \node[anchor=east, font=\scriptsize\bfseries\sffamily] at (-0.15,-1.00) {ASK};
  \draw[black!25] (0,-1.00) -- (9.10,-1.00);
  \draw[thick, netblue] plot[smooth] coordinates {(0.000,-1.000) (0.035,-1.000) (0.070,-1.000) (0.105,-1.000) (0.140,-1.000) (0.175,-1.000) (0.210,-1.000) (0.245,-1.000) (0.280,-1.000) (0.315,-1.000) (0.350,-1.000) (0.385,-1.000) (0.420,-1.000) (0.455,-1.000) (0.490,-1.000) (0.525,-1.000) (0.560,-1.000) (0.595,-1.000) (0.630,-1.000) (0.665,-1.000) (0.700,-1.000) (0.735,-0.824) (0.770,-0.715) (0.805,-0.715) (0.840,-0.824) (0.875,-1.000) (0.910,-1.176) (0.945,-1.285) (0.980,-1.285) (1.015,-1.176) (1.050,-1.000) (1.085,-0.824) (1.120,-0.715) (1.155,-0.715) (1.190,-0.824) (1.225,-1.000) (1.260,-1.176) (1.295,-1.285) (1.330,-1.285) (1.365,-1.176) (1.400,-1.000) (1.435,-1.000) (1.470,-1.000) (1.505,-1.000) (1.540,-1.000) (1.575,-1.000) (1.610,-1.000) (1.645,-1.000) (1.680,-1.000) (1.715,-1.000) (1.750,-1.000) (1.785,-1.000) (1.820,-1.000) (1.855,-1.000) (1.890,-1.000) (1.925,-1.000) (1.960,-1.000) (1.995,-1.000) (2.030,-1.000) (2.065,-1.000) (2.100,-1.000) (2.135,-0.824) (2.170,-0.715) (2.205,-0.715) (2.240,-0.824) (2.275,-1.000) (2.310,-1.176) (2.345,-1.285) (2.380,-1.285) (2.415,-1.176) (2.450,-1.000) (2.485,-0.824) (2.520,-0.715) (2.555,-0.715) (2.590,-0.824) (2.625,-1.000) (2.660,-1.176) (2.695,-1.285) (2.730,-1.285) (2.765,-1.176) (2.800,-1.000) (2.835,-0.824) (2.870,-0.715) (2.905,-0.715) (2.940,-0.824) (2.975,-1.000) (3.010,-1.176) (3.045,-1.285) (3.080,-1.285) (3.115,-1.176) (3.150,-1.000) (3.185,-0.824) (3.220,-0.715) (3.255,-0.715) (3.290,-0.824) (3.325,-1.000) (3.360,-1.176) (3.395,-1.285) (3.430,-1.285) (3.465,-1.176) (3.500,-1.000) (3.535,-1.000) (3.570,-1.000) (3.605,-1.000) (3.640,-1.000) (3.675,-1.000) (3.710,-1.000) (3.745,-1.000) (3.780,-1.000) (3.815,-1.000) (3.850,-1.000) (3.885,-1.000) (3.920,-1.000) (3.955,-1.000) (3.990,-1.000) (4.025,-1.000) (4.060,-1.000) (4.095,-1.000) (4.130,-1.000) (4.165,-1.000) (4.200,-1.000) (4.235,-1.000) (4.270,-1.000) (4.305,-1.000) (4.340,-1.000) (4.375,-1.000) (4.410,-1.000) (4.445,-1.000) (4.480,-1.000) (4.515,-1.000) (4.550,-1.000) (4.585,-1.000) (4.620,-1.000) (4.655,-1.000) (4.690,-1.000) (4.725,-1.000) (4.760,-1.000) (4.795,-1.000) (4.830,-1.000) (4.865,-1.000) (4.900,-1.000) (4.935,-0.824) (4.970,-0.715) (5.005,-0.715) (5.040,-0.824) (5.075,-1.000) (5.110,-1.176) (5.145,-1.285) (5.180,-1.285) (5.215,-1.176) (5.250,-1.000) (5.285,-0.824) (5.320,-0.715) (5.355,-0.715) (5.390,-0.824) (5.425,-1.000) (5.460,-1.176) (5.495,-1.285) (5.530,-1.285) (5.565,-1.176) (5.600,-1.000) (5.635,-1.000) (5.670,-1.000) (5.705,-1.000) (5.740,-1.000) (5.775,-1.000) (5.810,-1.000) (5.845,-1.000) (5.880,-1.000) (5.915,-1.000) (5.950,-1.000) (5.985,-1.000) (6.020,-1.000) (6.055,-1.000) (6.090,-1.000) (6.125,-1.000) (6.160,-1.000) (6.195,-1.000) (6.230,-1.000) (6.265,-1.000) (6.300,-1.000) (6.335,-1.000) (6.370,-1.000) (6.405,-1.000) (6.440,-1.000) (6.475,-1.000) (6.510,-1.000) (6.545,-1.000) (6.580,-1.000) (6.615,-1.000) (6.650,-1.000) (6.685,-1.000) (6.720,-1.000) (6.755,-1.000) (6.790,-1.000) (6.825,-1.000) (6.860,-1.000) (6.895,-1.000) (6.930,-1.000) (6.965,-1.000) (7.000,-1.000) (7.035,-0.824) (7.070,-0.715) (7.105,-0.715) (7.140,-0.824) (7.175,-1.000) (7.210,-1.176) (7.245,-1.285) (7.280,-1.285) (7.315,-1.176) (7.350,-1.000) (7.385,-0.824) (7.420,-0.715) (7.455,-0.715) (7.490,-0.824) (7.525,-1.000) (7.560,-1.176) (7.595,-1.285) (7.630,-1.285) (7.665,-1.176) (7.700,-1.000) (7.735,-1.000) (7.770,-1.000) (7.805,-1.000) (7.840,-1.000) (7.875,-1.000) (7.910,-1.000) (7.945,-1.000) (7.980,-1.000) (8.015,-1.000) (8.050,-1.000) (8.085,-1.000) (8.120,-1.000) (8.155,-1.000) (8.190,-1.000) (8.225,-1.000) (8.260,-1.000) (8.295,-1.000) (8.330,-1.000) (8.365,-1.000) (8.400,-1.000) (8.435,-1.000) (8.470,-1.000) (8.505,-1.000) (8.540,-1.000) (8.575,-1.000) (8.610,-1.000) (8.645,-1.000) (8.680,-1.000) (8.715,-1.000) (8.750,-1.000) (8.785,-1.000) (8.820,-1.000) (8.855,-1.000) (8.890,-1.000) (8.925,-1.000) (8.960,-1.000) (8.995,-1.000) (9.030,-1.000) (9.065,-1.000) (9.100,-1.000)};
  \node[anchor=east, font=\scriptsize\bfseries\sffamily] at (-0.15,-2.00) {FSK};
  \draw[black!25] (0,-2.00) -- (9.10,-2.00);
  \draw[thick, netblue] plot[smooth] coordinates {(0.000,-2.000) (0.035,-1.907) (0.070,-1.824) (0.105,-1.757) (0.140,-1.715) (0.175,-1.700) (0.210,-1.715) (0.245,-1.757) (0.280,-1.824) (0.315,-1.907) (0.350,-2.000) (0.385,-2.093) (0.420,-2.176) (0.455,-2.243) (0.490,-2.285) (0.525,-2.300) (0.560,-2.285) (0.595,-2.243) (0.630,-2.176) (0.665,-2.093) (0.700,-2.000) (0.735,-1.757) (0.770,-1.715) (0.805,-1.907) (0.840,-2.176) (0.875,-2.300) (0.910,-2.176) (0.945,-1.907) (0.980,-1.715) (1.015,-1.757) (1.050,-2.000) (1.085,-2.243) (1.120,-2.285) (1.155,-2.093) (1.190,-1.824) (1.225,-1.700) (1.260,-1.824) (1.295,-2.093) (1.330,-2.285) (1.365,-2.243) (1.400,-2.000) (1.435,-1.907) (1.470,-1.824) (1.505,-1.757) (1.540,-1.715) (1.575,-1.700) (1.610,-1.715) (1.645,-1.757) (1.680,-1.824) (1.715,-1.907) (1.750,-2.000) (1.785,-2.093) (1.820,-2.176) (1.855,-2.243) (1.890,-2.285) (1.925,-2.300) (1.960,-2.285) (1.995,-2.243) (2.030,-2.176) (2.065,-2.093) (2.100,-2.000) (2.135,-1.757) (2.170,-1.715) (2.205,-1.907) (2.240,-2.176) (2.275,-2.300) (2.310,-2.176) (2.345,-1.907) (2.380,-1.715) (2.415,-1.757) (2.450,-2.000) (2.485,-2.243) (2.520,-2.285) (2.555,-2.093) (2.590,-1.824) (2.625,-1.700) (2.660,-1.824) (2.695,-2.093) (2.730,-2.285) (2.765,-2.243) (2.800,-2.000) (2.835,-1.757) (2.870,-1.715) (2.905,-1.907) (2.940,-2.176) (2.975,-2.300) (3.010,-2.176) (3.045,-1.907) (3.080,-1.715) (3.115,-1.757) (3.150,-2.000) (3.185,-2.243) (3.220,-2.285) (3.255,-2.093) (3.290,-1.824) (3.325,-1.700) (3.360,-1.824) (3.395,-2.093) (3.430,-2.285) (3.465,-2.243) (3.500,-2.000) (3.535,-1.907) (3.570,-1.824) (3.605,-1.757) (3.640,-1.715) (3.675,-1.700) (3.710,-1.715) (3.745,-1.757) (3.780,-1.824) (3.815,-1.907) (3.850,-2.000) (3.885,-2.093) (3.920,-2.176) (3.955,-2.243) (3.990,-2.285) (4.025,-2.300) (4.060,-2.285) (4.095,-2.243) (4.130,-2.176) (4.165,-2.093) (4.200,-2.000) (4.235,-1.907) (4.270,-1.824) (4.305,-1.757) (4.340,-1.715) (4.375,-1.700) (4.410,-1.715) (4.445,-1.757) (4.480,-1.824) (4.515,-1.907) (4.550,-2.000) (4.585,-2.093) (4.620,-2.176) (4.655,-2.243) (4.690,-2.285) (4.725,-2.300) (4.760,-2.285) (4.795,-2.243) (4.830,-2.176) (4.865,-2.093) (4.900,-2.000) (4.935,-1.757) (4.970,-1.715) (5.005,-1.907) (5.040,-2.176) (5.075,-2.300) (5.110,-2.176) (5.145,-1.907) (5.180,-1.715) (5.215,-1.757) (5.250,-2.000) (5.285,-2.243) (5.320,-2.285) (5.355,-2.093) (5.390,-1.824) (5.425,-1.700) (5.460,-1.824) (5.495,-2.093) (5.530,-2.285) (5.565,-2.243) (5.600,-2.000) (5.635,-1.907) (5.670,-1.824) (5.705,-1.757) (5.740,-1.715) (5.775,-1.700) (5.810,-1.715) (5.845,-1.757) (5.880,-1.824) (5.915,-1.907) (5.950,-2.000) (5.985,-2.093) (6.020,-2.176) (6.055,-2.243) (6.090,-2.285) (6.125,-2.300) (6.160,-2.285) (6.195,-2.243) (6.230,-2.176) (6.265,-2.093) (6.300,-2.000) (6.335,-1.907) (6.370,-1.824) (6.405,-1.757) (6.440,-1.715) (6.475,-1.700) (6.510,-1.715) (6.545,-1.757) (6.580,-1.824) (6.615,-1.907) (6.650,-2.000) (6.685,-2.093) (6.720,-2.176) (6.755,-2.243) (6.790,-2.285) (6.825,-2.300) (6.860,-2.285) (6.895,-2.243) (6.930,-2.176) (6.965,-2.093) (7.000,-2.000) (7.035,-1.757) (7.070,-1.715) (7.105,-1.907) (7.140,-2.176) (7.175,-2.300) (7.210,-2.176) (7.245,-1.907) (7.280,-1.715) (7.315,-1.757) (7.350,-2.000) (7.385,-2.243) (7.420,-2.285) (7.455,-2.093) (7.490,-1.824) (7.525,-1.700) (7.560,-1.824) (7.595,-2.093) (7.630,-2.285) (7.665,-2.243) (7.700,-2.000) (7.735,-1.907) (7.770,-1.824) (7.805,-1.757) (7.840,-1.715) (7.875,-1.700) (7.910,-1.715) (7.945,-1.757) (7.980,-1.824) (8.015,-1.907) (8.050,-2.000) (8.085,-2.093) (8.120,-2.176) (8.155,-2.243) (8.190,-2.285) (8.225,-2.300) (8.260,-2.285) (8.295,-2.243) (8.330,-2.176) (8.365,-2.093) (8.400,-2.000) (8.435,-1.907) (8.470,-1.824) (8.505,-1.757) (8.540,-1.715) (8.575,-1.700) (8.610,-1.715) (8.645,-1.757) (8.680,-1.824) (8.715,-1.907) (8.750,-2.000) (8.785,-2.093) (8.820,-2.176) (8.855,-2.243) (8.890,-2.285) (8.925,-2.300) (8.960,-2.285) (8.995,-2.243) (9.030,-2.176) (9.065,-2.093) (9.100,-2.000)};
  \node[anchor=east, font=\scriptsize\bfseries\sffamily] at (-0.15,-3.00) {PSK};
  \draw[black!25] (0,-3.00) -- (9.10,-3.00);
  \draw[thick, netblue] plot[smooth] coordinates {(0.000,-3.000) (0.035,-2.824) (0.070,-2.715) (0.105,-2.715) (0.140,-2.824) (0.175,-3.000) (0.210,-3.176) (0.245,-3.285) (0.280,-3.285) (0.315,-3.176) (0.350,-3.000) (0.385,-2.824) (0.420,-2.715) (0.455,-2.715) (0.490,-2.824) (0.525,-3.000) (0.560,-3.176) (0.595,-3.285) (0.630,-3.285) (0.665,-3.176) (0.700,-3.000) (0.735,-3.176) (0.770,-3.285) (0.805,-3.285) (0.840,-3.176) (0.875,-3.000) (0.910,-2.824) (0.945,-2.715) (0.980,-2.715) (1.015,-2.824) (1.050,-3.000) (1.085,-3.176) (1.120,-3.285) (1.155,-3.285) (1.190,-3.176) (1.225,-3.000) (1.260,-2.824) (1.295,-2.715) (1.330,-2.715) (1.365,-2.824) (1.400,-3.000) (1.435,-2.824) (1.470,-2.715) (1.505,-2.715) (1.540,-2.824) (1.575,-3.000) (1.610,-3.176) (1.645,-3.285) (1.680,-3.285) (1.715,-3.176) (1.750,-3.000) (1.785,-2.824) (1.820,-2.715) (1.855,-2.715) (1.890,-2.824) (1.925,-3.000) (1.960,-3.176) (1.995,-3.285) (2.030,-3.285) (2.065,-3.176) (2.100,-3.000) (2.135,-3.176) (2.170,-3.285) (2.205,-3.285) (2.240,-3.176) (2.275,-3.000) (2.310,-2.824) (2.345,-2.715) (2.380,-2.715) (2.415,-2.824) (2.450,-3.000) (2.485,-3.176) (2.520,-3.285) (2.555,-3.285) (2.590,-3.176) (2.625,-3.000) (2.660,-2.824) (2.695,-2.715) (2.730,-2.715) (2.765,-2.824) (2.800,-3.000) (2.835,-3.176) (2.870,-3.285) (2.905,-3.285) (2.940,-3.176) (2.975,-3.000) (3.010,-2.824) (3.045,-2.715) (3.080,-2.715) (3.115,-2.824) (3.150,-3.000) (3.185,-3.176) (3.220,-3.285) (3.255,-3.285) (3.290,-3.176) (3.325,-3.000) (3.360,-2.824) (3.395,-2.715) (3.430,-2.715) (3.465,-2.824) (3.500,-3.000) (3.535,-2.824) (3.570,-2.715) (3.605,-2.715) (3.640,-2.824) (3.675,-3.000) (3.710,-3.176) (3.745,-3.285) (3.780,-3.285) (3.815,-3.176) (3.850,-3.000) (3.885,-2.824) (3.920,-2.715) (3.955,-2.715) (3.990,-2.824) (4.025,-3.000) (4.060,-3.176) (4.095,-3.285) (4.130,-3.285) (4.165,-3.176) (4.200,-3.000) (4.235,-2.824) (4.270,-2.715) (4.305,-2.715) (4.340,-2.824) (4.375,-3.000) (4.410,-3.176) (4.445,-3.285) (4.480,-3.285) (4.515,-3.176) (4.550,-3.000) (4.585,-2.824) (4.620,-2.715) (4.655,-2.715) (4.690,-2.824) (4.725,-3.000) (4.760,-3.176) (4.795,-3.285) (4.830,-3.285) (4.865,-3.176) (4.900,-3.000) (4.935,-3.176) (4.970,-3.285) (5.005,-3.285) (5.040,-3.176) (5.075,-3.000) (5.110,-2.824) (5.145,-2.715) (5.180,-2.715) (5.215,-2.824) (5.250,-3.000) (5.285,-3.176) (5.320,-3.285) (5.355,-3.285) (5.390,-3.176) (5.425,-3.000) (5.460,-2.824) (5.495,-2.715) (5.530,-2.715) (5.565,-2.824) (5.600,-3.000) (5.635,-2.824) (5.670,-2.715) (5.705,-2.715) (5.740,-2.824) (5.775,-3.000) (5.810,-3.176) (5.845,-3.285) (5.880,-3.285) (5.915,-3.176) (5.950,-3.000) (5.985,-2.824) (6.020,-2.715) (6.055,-2.715) (6.090,-2.824) (6.125,-3.000) (6.160,-3.176) (6.195,-3.285) (6.230,-3.285) (6.265,-3.176) (6.300,-3.000) (6.335,-2.824) (6.370,-2.715) (6.405,-2.715) (6.440,-2.824) (6.475,-3.000) (6.510,-3.176) (6.545,-3.285) (6.580,-3.285) (6.615,-3.176) (6.650,-3.000) (6.685,-2.824) (6.720,-2.715) (6.755,-2.715) (6.790,-2.824) (6.825,-3.000) (6.860,-3.176) (6.895,-3.285) (6.930,-3.285) (6.965,-3.176) (7.000,-3.000) (7.035,-3.176) (7.070,-3.285) (7.105,-3.285) (7.140,-3.176) (7.175,-3.000) (7.210,-2.824) (7.245,-2.715) (7.280,-2.715) (7.315,-2.824) (7.350,-3.000) (7.385,-3.176) (7.420,-3.285) (7.455,-3.285) (7.490,-3.176) (7.525,-3.000) (7.560,-2.824) (7.595,-2.715) (7.630,-2.715) (7.665,-2.824) (7.700,-3.000) (7.735,-2.824) (7.770,-2.715) (7.805,-2.715) (7.840,-2.824) (7.875,-3.000) (7.910,-3.176) (7.945,-3.285) (7.980,-3.285) (8.015,-3.176) (8.050,-3.000) (8.085,-2.824) (8.120,-2.715) (8.155,-2.715) (8.190,-2.824) (8.225,-3.000) (8.260,-3.176) (8.295,-3.285) (8.330,-3.285) (8.365,-3.176) (8.400,-3.000) (8.435,-2.824) (8.470,-2.715) (8.505,-2.715) (8.540,-2.824) (8.575,-3.000) (8.610,-3.176) (8.645,-3.285) (8.680,-3.285) (8.715,-3.176) (8.750,-3.000) (8.785,-2.824) (8.820,-2.715) (8.855,-2.715) (8.890,-2.824) (8.925,-3.000) (8.960,-3.176) (8.995,-3.285) (9.030,-3.285) (9.065,-3.176) (9.100,-3.000)};
\end{tikzpicture}
\caption{ASK: due ampiezze, qui portante o niente. FSK: due toni. PSK: la portante spostata di 0 o 180 gradi.}
\end{figure}

Quattro fasi portano 2 bit per simbolo: la \textbf{QPSK}. Frequenza e fase non possono variare indipendentemente, quindi gli schemi
reali combinano \textbf{ampiezza e fase}: la \textbf{QAM} (\emph{Quadrature Amplitude Modulation}).

\subsection{I diagrammi di costellazione}

\begin{figure}[H]
\centering
\begin{tikzpicture}[every node/.style={font=\scriptsize\sffamily}]
  \begin{scope}
    \draw[black!40, -{Stealth}] (-1.4,0) -- (1.4,0); \draw[black!40, -{Stealth}] (0,-1.4) -- (0,1.4);
    \foreach \x in {-0.7,0.7} \foreach \y in {-0.7,0.7} \fill[netblue] (\x,\y) circle (2.5pt);
    \draw[netred, thick, -{Stealth}] (0,0) -- (0.67,0.67); \node[text=netred, rotate=45] at (0.15,0.55) {ampiezza};
    \draw[netred] (0.4,0) arc (0:45:0.4); \node[text=netred] at (0.75,0.18) {fase};
    \node[font=\footnotesize\bfseries\sffamily] at (0,-1.75) {QPSK: 2 bit/simbolo};
  \end{scope}
  \begin{scope}[xshift=4.2cm]
    \draw[black!40, -{Stealth}] (-1.4,0) -- (1.4,0); \draw[black!40, -{Stealth}] (0,-1.4) -- (0,1.4);
    \foreach \x in {-1.05,-0.35,0.35,1.05} \foreach \y in {-1.05,-0.35,0.35,1.05} \fill[netblue] (\x,\y) circle (2.2pt);
    \node[font=\footnotesize\bfseries\sffamily] at (0,-1.75) {QAM-16: 4 bit/simbolo};
  \end{scope}
  \begin{scope}[xshift=8.4cm]
    \draw[black!40, -{Stealth}] (-1.4,0) -- (1.4,0); \draw[black!40, -{Stealth}] (0,-1.4) -- (0,1.4);
    \foreach \i in {0,...,7} \foreach \j in {0,...,7} \fill[netblue] (-1.05+\i*0.3,-1.05+\j*0.3) circle (1.6pt);
    \node[font=\footnotesize\bfseries\sffamily] at (0,-1.75) {QAM-64: 6 bit/simbolo};
  \end{scope}
\end{tikzpicture}
\caption{Ogni punto è un simbolo lecito: la distanza dall'origine è l'ampiezza, l'angolo è la fase.}
\end{figure}

\begin{itemize}
    \item Con $n$ punti si portano $\log_2 n$ bit per simbolo.
    \item Griglie quadrate, perché l'elettronica imposta un valore su ciascun asse più facilmente di un'ampiezza e di un angolo.
    \item Griglie più dense richiedono più $S/N$ per tenere separati i punti: è qui che si spende il budget di Shannon.
\end{itemize}

\subsection{Un punto sbagliato, quanti bit sbagliati?}

\begin{minipage}{0.4\linewidth}
\centering
\begin{tikzpicture}[every node/.style={font=\scriptsize\ttfamily}]
  \draw[black!40, -{Stealth}] (-2.1,0) -- (2.1,0) node[right, font=\scriptsize\sffamily] {I};
  \draw[black!40, -{Stealth}] (0,-2.1) -- (0,2.1) node[above, font=\scriptsize\sffamily] {Q};
  \foreach \r/\row in {0/{0000,0100,1100,1000}, 1/{0001,0101,1101,1001}, 2/{0011,0111,1111,1011}, 3/{0010,0110,1110,1010}} {
    \foreach \lab [count=\c from 0] in \row {
      \fill[netblue] (-1.5+\c,1.5-\r) circle (2pt);
      \node[below=5pt] at (-1.5+\c,1.5-\r) {\lab};
    }
  }
  \draw[netred, very thick] (0.5,0.5) circle (0.13);
  \foreach \x/\y in {0.5/1.5, 1.5/0.5, 0.5/-0.5, -0.5/0.5} \draw[netorange, thick] (\x,\y) circle (0.12);
\end{tikzpicture}
\end{minipage}\hfill
\begin{minipage}{0.56\linewidth}
Il rumore spinge il simbolo \texttt{1101} su un punto vicino. Quanti bit sono sbagliati? Numerati nell'ordine di conteggio, un vicino
di \texttt{0111} potrebbe essere \texttt{1000}: quattro bit sbagliati. Con un \textbf{codice di Gray} ogni vicino differisce in
\textbf{un solo bit}.
\begin{itemize}
    \item I vicini di \texttt{1101} sono \texttt{1100}, \texttt{1001}, \texttt{1111} e \texttt{0101}: un bit sbagliato ciascuno.
    \item Il rumore di solito sposta un simbolo solo di poco, quindi l'errore più comune costa un bit (che i codici del capitolo sugli
          errori possono poi correggere).
\end{itemize}
\end{minipage}

% ══════════════════════════════════════════════════════════════
\section{Il multiplexing}
% ══════════════════════════════════════════════════════════════

Una linea costa più o meno lo stesso qualunque sia la sua capacità: i soldi vanno nello scavo della trincea, non in ciò che ci si posa
dentro. Per questo una linea ad alta capacità è condivisa da molti segnali: è il \textbf{multiplexing}, per frequenza, per tempo, per
codice, e per lunghezza d'onda su una fibra.

\info{Multiplexing e accesso multiplo}{
    Abbiamo incontrato gli stessi nomi come modi in cui le stazioni condividono un canale broadcast: l'\textbf{accesso multiplo}. La
    differenza è \textbf{chi} fa la condivisione. Nel multiplexing un unico dispositivo, il \textbf{multiplexer}, combina i flussi a un
    capo e li separa all'altro. Nell'accesso multiplo le stazioni sono indipendenti e devono coordinarsi attraverso il canale stesso.
    Stessa divisione, problema diverso.
}

\subsection{Divisione di frequenza e bande di guardia}

\begin{figure}[H]
\centering
\begin{tikzpicture}[every node/.style={font=\scriptsize\sffamily}]
  \draw[-{Stealth}] (0,0) -- (2.6,0) node[right] {kHz};
  \fill[netblue!35] (0.3,0) -- (0.5,0.9) -- (1.9,0.9) -- (2.1,0) -- cycle; \node at (1.2,1.15) {una voce};
  \node[below] at (0.3,0) {0,3}; \node[below] at (2.1,0) {3,4};
  \draw[-{Stealth}, thick] (3.2,0.5) -- node[above, align=center] {spostata di\\60, 64, 68 kHz} (4.6,0.5);
  \begin{scope}[xshift=5cm]
    \draw[-{Stealth}] (0,0) -- (7.8,0) node[right] {kHz};
    \foreach \i/\c in {0/netblue!35, 1/netgreen!35, 2/netorange!40} {
      \fill[\c] (0.15+\i*2.4,0) -- (0.35+\i*2.4,0.9) -- (2.05+\i*2.4,0.9) -- (2.25+\i*2.4,0) -- cycle;
      \node at (1.2+\i*2.4,1.15) {canale \the\numexpr\i+1\relax};
    }
    \foreach \i/\k in {0/60, 1/64, 2/68, 3/72} \draw (\i*2.4,0.05) -- (\i*2.4,-0.1) node[below] {\k};
    \fill[netred!40] (2.25,0) rectangle (2.55,0.25);
    \draw[netred, -{Stealth}] (2.4,-0.75) node[below, text=netred] {banda di guardia} -- (2.4,0.1);
  \end{scope}
\end{tikzpicture}
\caption{FDM: ogni canale vocale, filtrato a 300--3400 Hz, viene spostato di una quantità diversa.}
\end{figure}

Ogni canale vocale viene alzato di una quantità diversa; nessuno si sovrappone più agli altri, e si possono sommare su una linea. Si
riservano 4000 Hz per canale: l'eccesso è una \textbf{banda di guardia}, perché i filtri reali non hanno bordi netti. La radio AM è la
stessa cosa nell'aria. Nel cuore della rete telefonica l'FDM ha ceduto il passo al TDM; sopravvive nella radio e nel cavo.

\subsection{OFDM}

\begin{figure}[H]
\centering
\begin{tikzpicture}[every node/.style={font=\scriptsize\sffamily}]
  \draw[very thick, netred, fill=netred!15] plot[smooth] coordinates {(-0.800,-0.146) (-0.720,-0.146) (-0.640,-0.142) (-0.560,-0.135) (-0.480,-0.125) (-0.400,-0.111) (-0.320,-0.094) (-0.240,-0.073) (-0.160,-0.051) (-0.080,-0.026) (0.000,0.000) (0.080,0.027) (0.160,0.054) (0.240,0.081) (0.320,0.107) (0.400,0.131) (0.480,0.153) (0.560,0.171) (0.640,0.186) (0.720,0.197) (0.800,0.204) (0.880,0.205) (0.960,0.202) (1.040,0.193) (1.120,0.179) (1.200,0.160) (1.280,0.136) (1.360,0.108) (1.440,0.075) (1.520,0.039) (1.600,-0.000) (1.680,-0.041) (1.760,-0.083) (1.840,-0.125) (1.920,-0.166) (2.000,-0.206) (2.080,-0.242) (2.160,-0.275) (2.240,-0.303) (2.320,-0.325) (2.400,-0.340) (2.480,-0.347) (2.560,-0.346) (2.640,-0.336) (2.720,-0.317) (2.800,-0.288) (2.880,-0.249) (2.960,-0.201) (3.040,-0.143) (3.120,-0.076) (3.200,0.000) (3.280,0.084) (3.360,0.175) (3.440,0.272) (3.520,0.374) (3.600,0.480) (3.680,0.589) (3.760,0.698) (3.840,0.807) (3.920,0.915) (4.000,1.019) (4.080,1.118) (4.160,1.211) (4.240,1.297) (4.320,1.373) (4.400,1.441) (4.480,1.497) (4.560,1.541) (4.640,1.574) (4.720,1.593) (4.800,1.600) (4.880,1.593) (4.960,1.574) (5.040,1.541) (5.120,1.497) (5.200,1.441) (5.280,1.373) (5.360,1.297) (5.440,1.211) (5.520,1.118) (5.600,1.019) (5.680,0.915) (5.760,0.807) (5.840,0.698) (5.920,0.589) (6.000,0.480) (6.080,0.374) (6.160,0.272) (6.240,0.175) (6.320,0.084) (6.400,0.000) (6.480,-0.076) (6.560,-0.143) (6.640,-0.201) (6.720,-0.249) (6.800,-0.288) (6.880,-0.317) (6.960,-0.336) (7.040,-0.346) (7.120,-0.347) (7.200,-0.340) (7.280,-0.325) (7.360,-0.303) (7.440,-0.275) (7.520,-0.242) (7.600,-0.206) (7.680,-0.166) (7.760,-0.125) (7.840,-0.083) (7.920,-0.041) (8.000,-0.000) (8.080,0.039) (8.160,0.075) (8.240,0.108) (8.320,0.136) (8.400,0.160) (8.480,0.179) (8.560,0.193) (8.640,0.202) (8.720,0.205) (8.800,0.204) (8.880,0.197) (8.960,0.186) (9.040,0.171) (9.120,0.153) (9.200,0.131) (9.280,0.107) (9.360,0.081) (9.440,0.054) (9.520,0.027) (9.600,0.000) (9.680,-0.026) (9.760,-0.051) (9.840,-0.073) (9.920,-0.094) (10.000,-0.111) (10.080,-0.125) (10.160,-0.135) (10.240,-0.142) (10.320,-0.146) (10.400,-0.146)};
  \draw[thick, netblue!70] plot[smooth] coordinates {(-0.800,-0.340) (-0.720,-0.347) (-0.640,-0.346) (-0.560,-0.336) (-0.480,-0.317) (-0.400,-0.288) (-0.320,-0.249) (-0.240,-0.201) (-0.160,-0.143) (-0.080,-0.076) (0.000,0.000) (0.080,0.084) (0.160,0.175) (0.240,0.272) (0.320,0.374) (0.400,0.480) (0.480,0.589) (0.560,0.698) (0.640,0.807) (0.720,0.915) (0.800,1.019) (0.880,1.118) (0.960,1.211) (1.040,1.297) (1.120,1.373) (1.200,1.441) (1.280,1.497) (1.360,1.541) (1.440,1.574) (1.520,1.593) (1.600,1.600) (1.680,1.593) (1.760,1.574) (1.840,1.541) (1.920,1.497) (2.000,1.441) (2.080,1.373) (2.160,1.297) (2.240,1.211) (2.320,1.118) (2.400,1.019) (2.480,0.915) (2.560,0.807) (2.640,0.698) (2.720,0.589) (2.800,0.480) (2.880,0.374) (2.960,0.272) (3.040,0.175) (3.120,0.084) (3.200,0.000) (3.280,-0.076) (3.360,-0.143) (3.440,-0.201) (3.520,-0.249) (3.600,-0.288) (3.680,-0.317) (3.760,-0.336) (3.840,-0.346) (3.920,-0.347) (4.000,-0.340) (4.080,-0.325) (4.160,-0.303) (4.240,-0.275) (4.320,-0.242) (4.400,-0.206) (4.480,-0.166) (4.560,-0.125) (4.640,-0.083) (4.720,-0.041) (4.800,-0.000) (4.880,0.039) (4.960,0.075) (5.040,0.108) (5.120,0.136) (5.200,0.160) (5.280,0.179) (5.360,0.193) (5.440,0.202) (5.520,0.205) (5.600,0.204) (5.680,0.197) (5.760,0.186) (5.840,0.171) (5.920,0.153) (6.000,0.131) (6.080,0.107) (6.160,0.081) (6.240,0.054) (6.320,0.027) (6.400,0.000) (6.480,-0.026) (6.560,-0.051) (6.640,-0.073) (6.720,-0.094) (6.800,-0.111) (6.880,-0.125) (6.960,-0.135) (7.040,-0.142) (7.120,-0.146) (7.200,-0.146) (7.280,-0.142) (7.360,-0.135) (7.440,-0.124) (7.520,-0.111) (7.600,-0.096) (7.680,-0.079) (7.760,-0.060) (7.840,-0.040) (7.920,-0.020) (8.000,-0.000) (8.080,0.020) (8.160,0.038) (8.240,0.056) (8.320,0.071) (8.400,0.085) (8.480,0.096) (8.560,0.104) (8.640,0.110) (8.720,0.113) (8.800,0.113) (8.880,0.111) (8.960,0.105) (9.040,0.098) (9.120,0.088) (9.200,0.076) (9.280,0.062) (9.360,0.048) (9.440,0.032) (9.520,0.016) (9.600,0.000) (9.680,-0.016) (9.760,-0.031) (9.840,-0.045) (9.920,-0.058) (10.000,-0.069) (10.080,-0.078) (10.160,-0.085) (10.240,-0.090) (10.320,-0.092) (10.400,-0.093)};
  \draw[thick, netblue!70] plot[smooth] coordinates {(-0.800,0.204) (-0.720,0.205) (-0.640,0.202) (-0.560,0.193) (-0.480,0.179) (-0.400,0.160) (-0.320,0.136) (-0.240,0.108) (-0.160,0.075) (-0.080,0.039) (0.000,-0.000) (0.080,-0.041) (0.160,-0.083) (0.240,-0.125) (0.320,-0.166) (0.400,-0.206) (0.480,-0.242) (0.560,-0.275) (0.640,-0.303) (0.720,-0.325) (0.800,-0.340) (0.880,-0.347) (0.960,-0.346) (1.040,-0.336) (1.120,-0.317) (1.200,-0.288) (1.280,-0.249) (1.360,-0.201) (1.440,-0.143) (1.520,-0.076) (1.600,0.000) (1.680,0.084) (1.760,0.175) (1.840,0.272) (1.920,0.374) (2.000,0.480) (2.080,0.589) (2.160,0.698) (2.240,0.807) (2.320,0.915) (2.400,1.019) (2.480,1.118) (2.560,1.211) (2.640,1.297) (2.720,1.373) (2.800,1.441) (2.880,1.497) (2.960,1.541) (3.040,1.574) (3.120,1.593) (3.200,1.600) (3.280,1.593) (3.360,1.574) (3.440,1.541) (3.520,1.497) (3.600,1.441) (3.680,1.373) (3.760,1.297) (3.840,1.211) (3.920,1.118) (4.000,1.019) (4.080,0.915) (4.160,0.807) (4.240,0.698) (4.320,0.589) (4.400,0.480) (4.480,0.374) (4.560,0.272) (4.640,0.175) (4.720,0.084) (4.800,0.000) (4.880,-0.076) (4.960,-0.143) (5.040,-0.201) (5.120,-0.249) (5.200,-0.288) (5.280,-0.317) (5.360,-0.336) (5.440,-0.346) (5.520,-0.347) (5.600,-0.340) (5.680,-0.325) (5.760,-0.303) (5.840,-0.275) (5.920,-0.242) (6.000,-0.206) (6.080,-0.166) (6.160,-0.125) (6.240,-0.083) (6.320,-0.041) (6.400,-0.000) (6.480,0.039) (6.560,0.075) (6.640,0.108) (6.720,0.136) (6.800,0.160) (6.880,0.179) (6.960,0.193) (7.040,0.202) (7.120,0.205) (7.200,0.204) (7.280,0.197) (7.360,0.186) (7.440,0.171) (7.520,0.153) (7.600,0.131) (7.680,0.107) (7.760,0.081) (7.840,0.054) (7.920,0.027) (8.000,0.000) (8.080,-0.026) (8.160,-0.051) (8.240,-0.073) (8.320,-0.094) (8.400,-0.111) (8.480,-0.125) (8.560,-0.135) (8.640,-0.142) (8.720,-0.146) (8.800,-0.146) (8.880,-0.142) (8.960,-0.135) (9.040,-0.124) (9.120,-0.111) (9.200,-0.096) (9.280,-0.079) (9.360,-0.060) (9.440,-0.040) (9.520,-0.020) (9.600,-0.000) (9.680,0.020) (9.760,0.038) (9.840,0.056) (9.920,0.071) (10.000,0.085) (10.080,0.096) (10.160,0.104) (10.240,0.110) (10.320,0.113) (10.400,0.113)};
  \draw[thick, netblue!70] plot[smooth] coordinates {(-0.800,0.113) (-0.720,0.113) (-0.640,0.110) (-0.560,0.104) (-0.480,0.096) (-0.400,0.085) (-0.320,0.071) (-0.240,0.056) (-0.160,0.038) (-0.080,0.020) (0.000,-0.000) (0.080,-0.020) (0.160,-0.040) (0.240,-0.060) (0.320,-0.079) (0.400,-0.096) (0.480,-0.111) (0.560,-0.124) (0.640,-0.135) (0.720,-0.142) (0.800,-0.146) (0.880,-0.146) (0.960,-0.142) (1.040,-0.135) (1.120,-0.125) (1.200,-0.111) (1.280,-0.094) (1.360,-0.073) (1.440,-0.051) (1.520,-0.026) (1.600,0.000) (1.680,0.027) (1.760,0.054) (1.840,0.081) (1.920,0.107) (2.000,0.131) (2.080,0.153) (2.160,0.171) (2.240,0.186) (2.320,0.197) (2.400,0.204) (2.480,0.205) (2.560,0.202) (2.640,0.193) (2.720,0.179) (2.800,0.160) (2.880,0.136) (2.960,0.108) (3.040,0.075) (3.120,0.039) (3.200,-0.000) (3.280,-0.041) (3.360,-0.083) (3.440,-0.125) (3.520,-0.166) (3.600,-0.206) (3.680,-0.242) (3.760,-0.275) (3.840,-0.303) (3.920,-0.325) (4.000,-0.340) (4.080,-0.347) (4.160,-0.346) (4.240,-0.336) (4.320,-0.317) (4.400,-0.288) (4.480,-0.249) (4.560,-0.201) (4.640,-0.143) (4.720,-0.076) (4.800,0.000) (4.880,0.084) (4.960,0.175) (5.040,0.272) (5.120,0.374) (5.200,0.480) (5.280,0.589) (5.360,0.698) (5.440,0.807) (5.520,0.915) (5.600,1.019) (5.680,1.118) (5.760,1.211) (5.840,1.297) (5.920,1.373) (6.000,1.441) (6.080,1.497) (6.160,1.541) (6.240,1.574) (6.320,1.593) (6.400,1.600) (6.480,1.593) (6.560,1.574) (6.640,1.541) (6.720,1.497) (6.800,1.441) (6.880,1.373) (6.960,1.297) (7.040,1.211) (7.120,1.118) (7.200,1.019) (7.280,0.915) (7.360,0.807) (7.440,0.698) (7.520,0.589) (7.600,0.480) (7.680,0.374) (7.760,0.272) (7.840,0.175) (7.920,0.084) (8.000,0.000) (8.080,-0.076) (8.160,-0.143) (8.240,-0.201) (8.320,-0.249) (8.400,-0.288) (8.480,-0.317) (8.560,-0.336) (8.640,-0.346) (8.720,-0.347) (8.800,-0.340) (8.880,-0.325) (8.960,-0.303) (9.040,-0.275) (9.120,-0.242) (9.200,-0.206) (9.280,-0.166) (9.360,-0.125) (9.440,-0.083) (9.520,-0.041) (9.600,-0.000) (9.680,0.039) (9.760,0.075) (9.840,0.108) (9.920,0.136) (10.000,0.160) (10.080,0.179) (10.160,0.193) (10.240,0.202) (10.320,0.205) (10.400,0.204)};
  \draw[thick, netblue!70] plot[smooth] coordinates {(-0.800,-0.093) (-0.720,-0.092) (-0.640,-0.090) (-0.560,-0.085) (-0.480,-0.078) (-0.400,-0.069) (-0.320,-0.058) (-0.240,-0.045) (-0.160,-0.031) (-0.080,-0.016) (0.000,0.000) (0.080,0.016) (0.160,0.032) (0.240,0.048) (0.320,0.062) (0.400,0.076) (0.480,0.088) (0.560,0.098) (0.640,0.105) (0.720,0.111) (0.800,0.113) (0.880,0.113) (0.960,0.110) (1.040,0.104) (1.120,0.096) (1.200,0.085) (1.280,0.071) (1.360,0.056) (1.440,0.038) (1.520,0.020) (1.600,-0.000) (1.680,-0.020) (1.760,-0.040) (1.840,-0.060) (1.920,-0.079) (2.000,-0.096) (2.080,-0.111) (2.160,-0.124) (2.240,-0.135) (2.320,-0.142) (2.400,-0.146) (2.480,-0.146) (2.560,-0.142) (2.640,-0.135) (2.720,-0.125) (2.800,-0.111) (2.880,-0.094) (2.960,-0.073) (3.040,-0.051) (3.120,-0.026) (3.200,0.000) (3.280,0.027) (3.360,0.054) (3.440,0.081) (3.520,0.107) (3.600,0.131) (3.680,0.153) (3.760,0.171) (3.840,0.186) (3.920,0.197) (4.000,0.204) (4.080,0.205) (4.160,0.202) (4.240,0.193) (4.320,0.179) (4.400,0.160) (4.480,0.136) (4.560,0.108) (4.640,0.075) (4.720,0.039) (4.800,-0.000) (4.880,-0.041) (4.960,-0.083) (5.040,-0.125) (5.120,-0.166) (5.200,-0.206) (5.280,-0.242) (5.360,-0.275) (5.440,-0.303) (5.520,-0.325) (5.600,-0.340) (5.680,-0.347) (5.760,-0.346) (5.840,-0.336) (5.920,-0.317) (6.000,-0.288) (6.080,-0.249) (6.160,-0.201) (6.240,-0.143) (6.320,-0.076) (6.400,0.000) (6.480,0.084) (6.560,0.175) (6.640,0.272) (6.720,0.374) (6.800,0.480) (6.880,0.589) (6.960,0.698) (7.040,0.807) (7.120,0.915) (7.200,1.019) (7.280,1.118) (7.360,1.211) (7.440,1.297) (7.520,1.373) (7.600,1.441) (7.680,1.497) (7.760,1.541) (7.840,1.574) (7.920,1.593) (8.000,1.600) (8.080,1.593) (8.160,1.574) (8.240,1.541) (8.320,1.497) (8.400,1.441) (8.480,1.373) (8.560,1.297) (8.640,1.211) (8.720,1.118) (8.800,1.019) (8.880,0.915) (8.960,0.807) (9.040,0.698) (9.120,0.589) (9.200,0.480) (9.280,0.374) (9.360,0.272) (9.440,0.175) (9.520,0.084) (9.600,0.000) (9.680,-0.076) (9.760,-0.143) (9.840,-0.201) (9.920,-0.249) (10.000,-0.288) (10.080,-0.317) (10.160,-0.336) (10.240,-0.346) (10.320,-0.347) (10.400,-0.340)};
  \draw[black!40, dashed] (1.60,-0.5) -- (1.60,1.75); \node[font=\scriptsize\sffamily] at (1.60,-0.75) {$f_1$};
  \draw[black!40, dashed] (3.20,-0.5) -- (3.20,1.75); \node[font=\scriptsize\sffamily] at (3.20,-0.75) {$f_2$};
  \draw[black!40, dashed] (4.80,-0.5) -- (4.80,1.75); \node[font=\scriptsize\sffamily] at (4.80,-0.75) {$f_3$};
  \draw[black!40, dashed] (6.40,-0.5) -- (6.40,1.75); \node[font=\scriptsize\sffamily] at (6.40,-0.75) {$f_4$};
  \draw[black!40, dashed] (8.00,-0.5) -- (8.00,1.75); \node[font=\scriptsize\sffamily] at (8.00,-0.75) {$f_5$};
  \draw[-{Stealth}] (-0.8,0) -- (11.2,0) node[right] {frequenza};
  \draw[{Stealth}-{Stealth}] (1.6,1.95) -- node[above] {separazione $\Delta f$} (3.2,1.95);
  \node[text=netred, anchor=west, text width=2.6cm] at (9.5,1.3) {in rosso una sottoportante: vale zero al centro di ogni altra};
\end{tikzpicture}
\caption{OFDM: le sottoportanti si sovrappongono, ma ciascuna è nulla al centro delle altre.}
\end{figure}

\begin{itemize}
    \item La banda viene tagliata in molte \textbf{sottoportanti}, ciascuna delle quali invia i propri dati con la QAM, accostate
          \textbf{senza bande di guardia}.
    \item Lo spettro di ciascuna sconfina nelle vicine, ma è zero al loro centro: campionate lì, le sottoportanti non si disturbano.
    \item Un flusso veloce diventa molti flussi lenti; una sottoportante colpita dal rumore riceve meno bit, o nessuno. Lo usano 802.11,
          il 4G, il cavo e l'ADSL.
\end{itemize}

\subsection{Divisione di tempo, fissa e statistica}

\begin{itemize}
    \item \textbf{TDM}: i flussi si alternano a turno (\emph{round-robin}), ciascuno padrone dell'intera banda per uno slot fisso. Gli
          estremi devono restare sincronizzati, e brevi tempi di guardia assorbono la deriva.
    \item \textbf{TDM statistico}: gli slot vanno a chi li chiede, non secondo un calendario: è la commutazione di pacchetto sotto un
          altro nome.
\end{itemize}

\begin{esempio}{La voce nel cuore della rete telefonica}
Un canale vocale da 4 kHz viene campionato 8000 volte al secondo (Nyquist: $2B$) con 8 bit per campione: \textbf{64 kb/s}. Una linea
\textbf{T1} intercala 24 di questi canali più un bit di trama: 193 bit ogni 125 $\mu$s, cioè \textbf{1,544 Mb/s}. Qui il
campionamento va nella direzione opposta a tutto il resto: un segnale analogico diventa bit, nel codec della centrale.
\end{esempio}

\subsection{Divisione di codice: quanto costano i chip}

\begin{itemize}
    \item Ogni tempo di bit è tagliato in $m$ \textbf{chip}. Una stazione invia il proprio codice di $m$ chip per un 1, e il suo negato
          per uno 0.
    \item Con codici ortogonali il ricevente moltiplica per il codice che vuole e ogni altra stazione sparisce: è l'algebra vista per
          l'accesso multiplo.
    \item Il costo: $m$ chip per bit richiedono $m$ volte la banda. 100 stazioni su 1 MHz: l'FDM dà a ciascuna 10 kHz; il CDMA dà a
          ciascuna l'intero MHz, a 100 chip per bit.
    \item I \textbf{codici di Walsh} forniscono $2^n$ codici ortogonali di lunghezza $2^n$. Restano ortogonali solo se i chip arrivano
          allineati nel tempo, cosa rara in una rete cellulare, che ha bisogno di altre soluzioni.
\end{itemize}

\subsection{Divisione di lunghezza d'onda}

\begin{itemize}
    \item Il \textbf{WDM} è la divisione di frequenza su una fibra: ogni segnale sulla propria lunghezza d'onda, combinati in un'unica
          fibra e separati all'altro capo da un \textbf{reticolo di diffrazione}, del tutto passivo e quindi molto affidabile.
    \item Perché conviene: un canale è limitato a pochi GHz dalla conversione elettro-ottica, in una banda di 25\,000 GHz. Le lunghezze
          d'onda parallele usano il resto.
    \item Nel 1990, 8 canali da 2,5 Gb/s; nel 2019, 160 canali e più di 16 Tb/s su una coppia di fibre.
    \item Gli amplificatori ottici rigenerano tutte le lunghezze d'onda insieme, senza riconvertirle una per una in elettricità.
\end{itemize}

% ══════════════════════════════════════════════════════════════
\section{Due casi: il modem e l'ADSL}
% ══════════════════════════════════════════════════════════════

\subsection{Perché il modem si è fermato a 33,6 kb/s}

Un modem telefonico invia i dati dove andrebbe la voce: 3100 Hz, a 2400 simboli al secondo. Generazione dopo generazione, costellazioni
più dense: 4, 6, 12, 14 bit per simbolo, cioè 9600, 14\,400, 28\,800, 33\,600 b/s.

\warning{Perché la progressione si è fermata a 14 bit per simbolo?}{
    Per \textbf{Shannon}. Attraverso due doppini d'utente medi, la capacità del canale vocale è di circa 35 kb/s: per andare più
    veloci serve un'altra fisica, o altri fili. Il trucco dei 56k: il provider si collega in digitale, eliminando uno dei due tratti
    analogici rumorosi. Il limite diventa la centrale: 8000 campioni da 7 bit utilizzabili, 56 kb/s.
}

\subsection{ADSL: togliere il filtro, dividere la linea}

\begin{figure}[H]
\centering
\begin{tikzpicture}[every node/.style={font=\scriptsize\sffamily}]
  \draw[-{Stealth}] (0,0) -- (12,0) node[right] {kHz};
  \fill[netgreen!40] (0,0) rectangle (0.6,0.8); \node at (0.3,1.05) {voce};
  \fill[netorange!40] (1.2,0) rectangle (3,0.8); \node at (2.1,1.05) {upstream};
  \fill[netblue!35] (3.1,0) rectangle (11.4,1.3); \node at (7.25,1.55) {downstream};
  \foreach \x in {3.4,3.7,...,11.2} \draw[white, thin] (\x,0) -- (\x,1.3);
  \foreach \x/\k in {0/0, 1.2/25, 3.05/138, 11.4/1104} \draw (\x,0) -- (\x,-0.1) node[below] {\k};
  \node[anchor=east, text=black!60] at (12,-0.6) {non in scala};
\end{tikzpicture}
\caption{ADSL: voce in basso, poi un canale di upstream e uno più grande di downstream, divisi in sottocanali.}
\end{figure}

\begin{itemize}
    \item I 3100 Hz erano un \textbf{filtro} in centrale, non un limite del filo: aggirato il filtro, il doppino offre circa 1,1 MHz.
    \item L'OFDM, qui chiamato \textbf{DMT} (\emph{Discrete MultiTone}), lo taglia in 256 canali: voce, alcuni vuoti, l'upstream e un
          downstream più grande. È la A (\emph{asymmetric}) di ADSL.
    \item Ogni canale usa la QAM a circa 4000 simboli al secondo, da 15 bit per simbolo dove $S/N$ è alto fino a nessuno dove non lo è.
\end{itemize}

Doppino, banda passante, QAM, OFDM, Shannon: tutto ciò che abbiamo visto, sul filo che entra in casa.

% ══════════════════════════════════════════════════════════════
\section{Il proprio collegamento, visto dall'adattatore}
% ══════════════════════════════════════════════════════════════

\begin{lstlisting}[language=bash]
$ ethtool enp3s0 | grep -E 'Speed|Duplex|Port'
    Speed: 1000Mb/s
    Duplex: Full
    Port: Twisted Pair
$ iw dev wlp2s0 link | grep -E 'freq|signal|tx bitrate'
    freq: 5180
    signal: -52 dBm
    tx bitrate: 866.7 MBit/s VHT-MCS 9 80MHz short GI VHT-NSS 2
\end{lstlisting}

\begin{itemize}
    \item Cablato: un gigabit su quattro doppini, full-duplex, negoziato dai due adattatori.
    \item \texttt{80MHz}: la banda analogica, la $B$ di Shannon. \texttt{-52 dBm}: la potenza ricevuta, circa 6 nanowatt.
    \item \texttt{MCS 9}: 256-QAM su sottoportanti OFDM, con due flussi spaziali (\texttt{NSS 2}, due antenne).
    \item Allontanandosi dall'access point si vede l'MCS scendere: l'adattatore sceglie la costellazione più densa che $S/N$ consente.
\end{itemize}

\begin{esercizio}{Nyquist, Shannon e codici}
\begin{enumerate}[label=(\alph*)]
    \item Un canale senza rumore da 3 kHz porta un segnale a 4 livelli. Qual è la velocità massima?
    \item Lo stesso canale ha in realtà $S/N = 30$ dB. Qual è la sua capacità? Quanti livelli servirebbero, secondo Nyquist, per
          raggiungerla?
    \item Quanti baud servono per trasmettere 100 Mb/s con NRZ, con Manchester e con 4B/5B?
    \item Verificare la velocità di una linea T1 e calcolare quanta capacità va al solo bit di trama.
    \item Il Wi-Fi dell'ultimo esempio usa 80 MHz. Con $S/N = 30$ dB, qual è il limite di Shannon per un singolo flusso? Come può
          l'adattatore dichiarare 866,7 Mb/s?
\end{enumerate}

\svolgimento
\begin{enumerate}[label=(\alph*)]
    \item $R_{\max} = 2 \cdot 3000 \cdot \log_2 4 = 12\,000$ b/s.
    \item $30$ dB significa $S/N = 1000$: $C = 3000 \cdot \log_2 1001 \approx 3000 \cdot 9{,}97 \approx 29{,}9$ kb/s. Livelli:
          $V = \sqrt{1001} \approx 32$, cioè 5 bit per simbolo a 6000 simboli al secondo.
    \item NRZ: un simbolo per bit, 100 Mbaud. Manchester: due transizioni per bit, 200 Mbaud. 4B/5B: 5 simboli ogni 4 bit,
          125 Mbaud (è il valore di 100BASE-TX).
    \item $24 \cdot 8 + 1 = 193$ bit ogni 125 $\mu$s, cioè $193 \cdot 8000 = 1\,544\,000$ b/s. Il bit di trama vale $8000$ b/s,
          circa lo 0,5\%.
    \item $C = 80 \cdot 10^6 \cdot \log_2 1001 \approx 797$ Mb/s per flusso. Con due flussi spaziali (due antenne, due percorsi
          indipendenti) il limite si applica a ciascun flusso: 866,7 Mb/s in totale sono circa 433 per flusso, ben sotto il tetto.
\end{enumerate}
\end{esercizio}

\Needspace{20\baselineskip}
% ══════════════════════════════════════════════════════════════
\section{Riepilogo}
% ══════════════════════════════════════════════════════════════

\riepilogo{
\begin{itemize}
    \item I bit diventano segnali (tensione, luce, radio) su mezzi \textbf{guidati} (doppino, coassiale, fibra) o \textbf{non guidati}
          (radio), ciascuno con la sua banda, la sua portata e il suo punto debole.
    \item Un canale lascia passare una banda di frequenze. \textbf{Nyquist} limita quanto in fretta si può segnalare ($2B\log_2 V$),
          \textbf{Shannon} quanti bit si possono portare: $B \log_2(1 + S/N)$.
    \item \textbf{Banda base}: i codici di linea scambiano banda con recupero dell'orologio e bilanciamento: Manchester, NRZI, 4B/5B,
          scrambling, 8B/10B.
    \item \textbf{Banda passante}: ampiezza, frequenza e fase di una portante; costellazioni dense quanto il rumore consente, con
          codice di Gray.
    \item Una linea, molti segnali: per frequenza, tempo, codice e lunghezza d'onda, e l'\textbf{OFDM}, che accosta sottoportanti senza
          bande di guardia.
    \item La pila è completa, dall'applicazione al filo. Prossimo passo: che cosa succede quando qualcuno su di essa è ostile.
\end{itemize}
}
